\SGAchapter{XIII}{问题与猜想}{Problèmes et conjectures}
\SourcePage{143}
\SGAsection{1}{整体结果与局部结果的关系；与对偶性有关的仿射问题}
众所周知，许多有关\Term{projective-scheme}{射影概形} $X$ 的命题，都可以表述为有关某个\Term{graded-ring}{分次环}，或者更好地说，某个\Term{complete}{完备}\Term{local-ring}{局部环}的命题：这个环就是 $X$ 的\Term{homogeneous-coordinate-ring}{齐次坐标环}（即 $X$ 的\Term{projecting-cone}{投影锥} $\widetilde X$ 的\Term{affine-ring}{仿射环}），或其\Term{completion}{完备化}（即 $\widetilde X$ 的顶点处\Term{local-ring}{局部环}的\Term{completion}{完备化}）。这样重新表述的意义在于，常常可以从已知的整体结果出发，猜测乃至证明关于\Term{complete}{完备}\Term{noetherian}{诺特}\Term{local-ring}{局部环}的类似结果，而这些环比整体命题中实际出现的环更一般，例如不必具有\Term{equal-characteristic}{等特征}。因此，\Term{projective-space}{射影空间}的 Serre \Term{duality}{对偶}定理（第 XII 讲 1.1）启发了有用的\Term{local-duality}{局部对偶}定理（第 V 讲 2.1）。Serre 关于\Term{projective-space}{射影空间}上\Term{coherent}{凝聚}模\Term{cohomology}{上同调}的基本定理（\Term{finite}{有限}性以及当 $n$ 很大时 $H^i(X,F(n))$ 的渐近性态，参见 EGA III 2.2.1），推广为\Term{local-scheme}{局部}\Term{invariants}{不变量} $H^i_{\mathfrak m}(M)$ 的结构定理，见第 V 讲第 3 节。同样，关于\Term{fundamental-group}{基本群}与 \Term{picard-group}{Picard 群}的 Lefschetz 定理（“\Term{equivalence}{等价}判据”），在经典情形中早已为人熟知，后来又推广到任意\Term{base-field}{基域}，它们启发了第 X、XI 讲的“\Term{local-scheme}{局部}”Lefschetz 定理。当然，\Term{local-scheme}{局部}定理反过来又是获得整体命题的宝贵工具。例如，\Term{local-duality}{局部对偶}性使我们能够用某些\Term{local-scheme}{局部}\Term{invariants}{不变量} $H^i(F_x)$ 的消失来表述整体渐近性质 XII 1.3 (i)。更实质地说，\Term{local-scheme}{局部} Lefschetz 定理蕴含某些\Term{local-complete-intersection}{局部完全交}环的“\Term{purity}{纯性}”或\Term{parafactorial}{副阶乘}性（X 3.4 和 XI 3.13），因而能够在整体 Lefschetz 定理中去掉某些\Term{nonsingular}{非奇异}性假设，如 X 3.5、3.6、3.7。

\SourcePage{144}
有关域 $k$ 上\Term{projective-scheme}{射影概形}的定理，还有另一种有用的推广，即把 $k$ 换成一般\Term{base-scheme}{基概形}。这样，EGA III 的后续部分将给出\footnote{编者注（1）：事实上，该书并未包含这项推广；见下文注释及原书第 2 页编者注（4）。}Serre \Term{duality}{对偶}性在这个方向的推广\footnote{参见第 IV 讲末尾引用的 Hartshorne 讨论班。}；$H^i(X,F(n))$ 的\Term{finite}{有限}性与渐近性态定理已在 EGA III 2.2.1 中对一般\Term{base-scheme}{基概形}陈述；最后，如 XII 4.9 所示，借助\Term{local-scheme}{局部}定理 XII 4.7，也能对\Term{projective-morphism}{射影态射}建立 Lefschetz 定理。当然，在一般\Term{base-scheme}{基概形}上工作，也会导向本质上全新的命题，如“\Term{comparison-theorem}{比较定理}”EGA III 4.15 和\Term{sheaf}{层}的\Term{existence-theorem}{存在定理} EGA III 5.1.4（我们也已在第 IX 讲看到，它们与关于 $\pi_1$ 和 $\operatorname{Pic}$ 的 Lefschetz 定理，归根结底依赖于同样的\Term{cohomology}{上同调}性质的关键定理）。

因此，必须找出能够同时涵盖刚才所说的、有关域上\Term{projective-scheme}{射影概形}的命题的两种推广的定理。适于这种共同推广的自然\Term{object}{对象}，是对某个 $I$ \Term{adic-topology}{进拓扑}\Term{separated}{分离}而\Term{complete}{完备}的\Term{noetherian}{诺特}环。从这个观点对它们的研究，还没有认真展开；在我看来，这在目前是\Term{coherent}{凝聚}\Term{sheaf}{层}\Term{local-scheme}{局部}理论中最有意思的课题。下面是这一方向的一个典型问题。

\begin{SGAresult}{猜想}{1.1（“第二\Term{affine}{仿射}\Term{finiteness-theorem}{有限性定理}”）}
\footnote{本猜想及下面的猜想 1.2 都是错误的，R. Hartshorne 已指出这一点，见《Affine duality and cofinite modules》，Invent. Math. 9 (1969/70)，145--164 页，第 3 节。}\footnote{编者注（2）：不过，若 $A$ 是\Term{complete}{完备}\Term{local-ring}{局部环}（分别地，是正\Term{characteristic}{特征}\Term{regular}{正则}环），且 $J$ 是\Term{maximal-ideal}{极大理想}，那么对\Term{finitely-generated}{有限生成}的 $M$（分别地，对 $M=A$），命题为真；参见 Hartshorne R., “Affine duality and cofinite modules”, Invent. Math. 9 (1969/70), 145--164 页，推论 1.4；分别参见 Huneke C. \& Sharp R., “Bass numbers of local cohomology modules”, Trans. Amer. Math. Soc. 339 (1993), 第 2 期，765--779 页，其中还包含远为更强的结果。关于处理零\Term{characteristic}{特征}的完全不同的方法（$D$ 模），见 Lyubeznik G., “Finiteness properties of local cohomology modules (an application of D-modules to commutative algebra)”, Invent. Math. 113 (1993), 第 1 期，41--55 页；同一作者的 “Finiteness properties of local cohomology modules for regular local rings of mixed characteristic: the unramified case”, Comm. Algebra 28 (2000), 第 12 期，5867--5882 页，Robin Hartshorne 纪念专辑；以及 “Finiteness properties of local cohomology modules: a characteristic-free approach”, J. Pure Appl. Algebra 151 (2000), 第 1 期，43--50 页。在 Hartshorne 的推动下，\Term{cofinite}{余有限}模的概念后来发生了变化：若 $M$ 的\Term{support}{支撑}包含于 $V(J)$，且所有 $\operatorname{Ext}^i_A(A/J,H^i_J(M))$ 都\Term{finitely-generated}{有限生成}，就称 $M$ 为 $J$ \Term{cofinite}{余有限}的。关于这一课题，例如参见 Delfino D. \& Marley Th., “Cofinite modules and local cohomology”, J. Pure Appl. Algebra 121 (1997), 第 1 期，45--52 页。}
设 $M$ 是\Term{noetherian}{诺特}环 $A$ 上的\Term{finitely-generated}{有限生成}模（必要时假设 $A$ 是某个\Term{regular}{正则}环的商），$J$ 是 $A$ 的\Term{ideal}{理想}。证明模 $H^i_J(M)$ 是“$J$ \Term{cofinite}{余有限}的”，即模
\SourcePage{145}
\[
\operatorname{Hom}_A(A/J,H^i_J(M))
\]
\Term{finitely-generated}{有限生成}。
\end{SGAresult}

回忆，$H^i_J(M)$ 表示第 I 讲中的模 $H^i_Y(X,\widetilde M)$（其中 $X=\operatorname{Spec}(A)$，$Y=V(J)$）；第 II 讲把它解释为 \Term{koszul-complex}{Koszul 复形}\Term{cohomology}{上同调}的\Term{direct-limit}{归纳极限}，而在 $i\geq2$ 时，它也就是模 $H^{i-1}(X-Y,\widetilde M)$。实际上，1.1 应当来自一个更精确的命题：$H^i_J(M)$ 属于\Term{support}{支撑}包含于 $Y=V(J)$ 的 $A$ 模\Term{category}{范畴} $\underline C_J$ 的某个合适的\Term{abelian-subcategory}{阿贝尔子范畴} $\underline D_J$，并且 $H\in\operatorname{Ob}\underline D_J$ 蕴含 $H$ 是 $J$ \Term{cofinite}{余有限}的。（注意：\Term{support}{支撑}包含于 $V(J)$ 的 $J$ \Term{cofinite}{余有限}模 $H$ 所成的\Term{category}{范畴}，很遗憾并不对取商\Term{stable}{稳定}！）于是，核心问题在于定义 $\underline D_J$。更精确地说，问题 1.1 的解答（至少在 $A$ 是某个\Term{regular}{正则}环的商时）应当来自一种\Term{duality}{对偶}理论，同时推广\Term{local-duality}{局部对偶}性以及前面提到的\Term{projective-morphism}{射影态射}\Term{duality}{对偶}理论，其形式应如下。

\begin{SGAresult}{猜想}{1.2（“\Term{affine}{仿射}\Term{duality}{对偶}性”）}
\footnote{这个猜想按原来的表述是错误的，不过 R. Hartshorne 已以相当接近的形式证明了它，见 “Affine duality and cofinite modules”, Invent. Math. 9 (1969/70)，145--164 页。}
假设 $A$ \Term{regular}{正则}，并且对 $J$ \Term{adic-topology}{进拓扑}\Term{separated}{分离}而\Term{complete}{完备}。设 $C^\bullet(A)$ 是 $A$ 的一个\Term{injective-resolution}{内射消解}。

(i) 证明\Term{functor}{函子}
\[
D_J:L_\bullet\longmapsto\operatorname{Hom}_J(L_\bullet,C^\bullet(A))
\]
是\Term{fully-faithful}{充分忠实}的。其源\Term{category}{范畴}为每一维均\Term{free}{自由}且\Term{finitely-generated}{有限生成}、次数上有界的 $A$ 模\Term{complex}{复形}\Term{category}{范畴}（\Term{morphism}{态射}为模去\Term{homotopy}{同伦}的\Term{complex}{复形}\Term{homomorphism}{同态}）；其目标\Term{category}{范畴}为每一维均\Term{injective}{内射}、次数上有界的 $A$ 模\Term{complex}{复形} $K^\bullet$ 的\Term{category}{范畴}（\Term{homomorphism}{同态}按相同方式定义）。

(ii) 证明，对于每个形如 $D_J(L_\bullet)$ 的 $K^\bullet$，$H^i(K^\bullet)$（即 $\operatorname{Ext}^i_Y(X;L_\bullet,\mathcal O_X)$）都是 $J$ \Term{cofinite}{余有限}的。

(iii) 更精确地，证明与某个形如 $D_J(L_\bullet)$ 的\Term{complex}{复形}\Term{homotopy}{同伦}的 $K^\bullet$，可以用 $H^i(K^\bullet)$ 的\Term{finite}{有限}性性质来刻画，这些性质比 (ii) 中所考虑的更强，例如要求 $H^i(K^\bullet)\in\operatorname{Ob}\underline D_J$，其中 $\underline D_J$ 是如上所述的某个合适的\Term{abelian-category}{阿贝尔范畴}。
\end{SGAresult}

注意，当 $A$ 是\Term{local-ring}{局部环}而 $J$ 是它的\Term{ideal-of-definition}{定义理想}时（参见第 IV 讲），以及当 $J$ 是零\Term{ideal}{理想}时，这个问题已得到肯定解答。在这两个特殊情形，可以直接把 $\underline D_J$ 取为\Term{support}{支撑}在 $V(J)$ 上的 $J$ \Term{cofinite}{余有限}模\Term{category}{范畴}（在第二个情形，这只是说取 $A$ 上\Term{finitely-generated}{有限生成}模的\Term{category}{范畴}）。
\SourcePage{146}
若猜想 1.2 在一般情形下有肯定解答，把 $L_\bullet$ 取为 $M$ 的\Term{finitely-generated}{有限生成}\Term{free-resolution}{自由消解}的\Term{duality}{对偶}，就会得到 1.1 的肯定解答。另一方面，1.1 的肯定解答又会肯定地回答下列猜想的第一部分；我们把它表述为“整体”形式。

\begin{SGAresult}{猜想}{1.3}
设 $X\subset\mathbf P^r_k$ 是标准\Term{projective-scheme}{射影概形}的一个\Term{closed-subscheme}{闭子概形}，\Term{local-scheme}{局部}是\Term{complete-intersection}{完全交}，且每个\Term{irreducible-component}{不可约分支}的\Term{codimension}{余维数}都 $\geq s$。令 $U=\mathbf P^r_k-X$。

(i) 证明，对于 $U$ 上任意\Term{coherent}{凝聚}模 $F$，都有
\[
\dim H^i(U,F)<+\infty\qquad\text{当 }i\geq s.
\]
\footnote{当 $X$ \Term{smooth}{光滑}时，R. Hartshorne 已证明本猜想的 (i)，参见 Ample subvarieties of algebraic varieties，与 C. Musili 合作撰写的笔记，Lect. Notes in Math., 第 156 卷，Springer-Verlag, Berlin--New York, 1970，定理 III.5.2。此作者也找到了 (ii) 的例子，见 “Cohomological dimension of algebraic varieties”, Ann. of Math. (2) 88 (1968)，403--450 页，449 页的例子。}

(ii) 给出一个 $X$ \Term{connected}{连通}且\Term{regular}{正则}、而
\[
H^s(U,F)\ne0
\]
的例子。
\end{SGAresult}
要看出 (i) 是 1.1 的特例，把
\[
H^i(U,F(\mathord\cdot))=\bigoplus_n H^i(U,F(n))
 =H^i(\mathbf E^{r+1}-\widetilde X,\widetilde F)
\]
视为 $\mathbf P^r$ 的\Term{projecting-cone}{投影锥} $\mathbf E^{r+1}$ 的\Term{affine-ring}{仿射环} $k[t_0,\ldots,t_r]$ 上的模。这个模正是 $H^{i+1}_J(M)$，其中 $J$ 是 $X$ 的\Term{projecting-cone}{投影锥} $\widetilde X$ 在 $\mathbf E^{r+1}$ 中的\Term{ideal}{理想}。另一方面，对 $X$ 的假设蕴含：在 $\mathbf E^{r+1}$ 的每个异于原点的点处，$\widetilde X$ 也是\Term{codimension}{余维数} $\geq s$ 的\Term{complete-intersection}{完全交}；因而当 $i\geq s$ 时，$H^{i+1}_J(M)$ 在原点以外为零。所以，若如 1.1 所要求的，它是 $J$ \Term{cofinite}{余有限}的，则更是 $\mathfrak m$ \Term{cofinite}{余有限}的，而这容易蕴含其每个次数的部分都\Term{finite}{有限}维。\footnote{编者注（3）：Hartshorne 在 “Cohomological dimension of algebraic varieties”, Ann. of Math. (2) 88 (1968), 403--450 页证明：当 $F$ \Term{coherent}{凝聚}且 $X$ \Term{dimension}{维数}为正时，\Term{cohomology}{上同调} $H^{n-1}(\mathbf P^n_k-X,F)$ 为零（$k$ \Term{algebraically-closed}{代数闭}）。实际上，主要借助 Serre \Term{duality}{对偶}性和 Lichtenbaum 定理——\Term{irreducible}{不可约}而非\Term{complete}{完备}的\Term{quasi-projective-morphism}{拟射影}簇，其最高维的\Term{coherent}{凝聚}\Term{sheaf}{层}\Term{cohomology}{上同调}消失——可归结为证明形式\Term{completion}{完备化} $\widehat{\mathbf P}{}^n_k$ 与 $X$ 具有相同的\Term{rational-function-field}{有理函数域}。这是困难的一点（所引文献定理 7.2）；用 Hironaka H. \& Matsumura H., “Formal functions and formal embeddings”, J. Math. Soc. Japan 20 (1968), 52--82 页中的术语来说，$\mathbf P^n_k$ 是 G3 的。这两位作者独立证明了上述结果，事实上还得到了强得多的结果。他们证明了 $X$ \Term{universal}{普遍}为 G3，并计算了\Term{abelian-variety}{阿贝尔簇}沿正维子簇的形式\Term{completion}{完备化}的\Term{rational-function-field}{有理函数域}。如今经典的 G1、G2、G3 条件，首次出现在这篇文章中。}还应注意，对于 $\mathbf P^3$ 中一条\Term{irreducible}{不可约}\Term{nonsingular}{非奇异}曲线 $X$，就已经可以提出猜想 1.3；在这种情形下，我们不知道 $H^2(\mathbf P^3-X,\mathcal O_X(n))$
\SourcePage{147}
是否\Term{finite}{有限}维，也不知道它们是否必然为零。\footnote{R. Hartshorne（“Ample vector bundles”, Publ. Math. Inst. Hautes Études Sci. 29 (1966), 63--94 页，定理 8.1）和 H. Hironaka 刚刚给出了肯定解答。}我们甚至不知道，是否存在 $\mathbf P^3$ 中的\Term{irreducible}{不可约}曲线，不能在\Term{set}{集合}意义下表示为两个\Term{hypersurface}{超曲面}的交。\footnote{编者注（4）：见猜想 3.5 及其相应注释。}

\begin{SGAresult}{问题}{1.4}
给出“\Term{comparison-theorem}{比较定理}”EGA III 4.1.5 的\Term{affine}{仿射}变体，把它表述为\Term{functor}{函子} $H^i_J$ 与某些\Term{inverse-limit}{逆极限}交换的定理。
\end{SGAresult}

最后，在这一思路下，我曾提出下述问题：设 $A$ 是\Term{complete}{完备}\Term{noetherian}{诺特}\Term{regular}{正则}\Term{local-ring}{局部环}，$K$ 是其\Term{fraction-field}{分式域}，证明对所有 $i$ 都有 $\operatorname{Ext}^i_A(K,A)=0$。M. Auslander 当即给出肯定解答；\Term{regularity}{正则性}假设可换成 $A$ 为\Term{integral-domain}{整环}，并且实际上，只要 $M$ 是 $A$ 上\Term{finitely-generated}{有限生成}模，就对所有 $i$ 都有 $\operatorname{Ext}^i_A(K,M)=0$。这直接来自 Auslander 的下述命题：若 $A$ 是\Term{complete}{完备}\Term{noetherian}{诺特}\Term{local-ring}{局部环}，则对于 $A$ 上任意\Term{finitely-generated}{有限生成}模 $M$，\Term{functor}{函子} $\operatorname{Ext}^i_A(\mathord\cdot,M)$ 把\Term{direct-limit}{归纳极限}变为\Term{inverse-limit}{逆极限}。\footnote{编者注（5）：写成 $K=\varinjlim_{a\ne0}A[1/a]$，并注意 $a\ne0$ 是 $A$ \Term{regular}{正则}的。}

\SGAsection{2}{与 $\pi_0$ 有关的问题：局部 Bertini 定理}
设 $A$ 是\Term{complete}{完备}\Term{noetherian}{诺特}\Term{local-ring}{局部环}，$f$ 是其\Term{maximal-ideal}{极大理想}中的元素，$X=\operatorname{Spec}(A)$，$Y=\operatorname{Spec}(A/fA)$。使用\Term{local-scheme}{局部}“Lefschetz”技巧，可以给出以 $X'=X-\{\mathfrak m\}$ 的假设表述的、使 $Y'=X'\cap Y$ \Term{connected}{连通}的判据。例如，下述条件就足够：a) $X'$ \Term{connected}{连通}；b) 对 $X'$ 的每个\Term{closed-point}{闭点} $x$ 都有 $\operatorname{prof}\mathcal O_{X',x}\geq2$；c) $f$ 是 $A$ \Term{regular}{正则}的。不过，应注意假设 b) 和 c) 并非纯\Term{topology}{拓扑}性质，例如，把 $X$ 换成 $X_{\mathrm{red}}$ 时它们并不保持不变。对于域上\Term{projective-scheme}{射影概形} $X'$ 与其\Term{hyperplane-section}{超平面截面} $Y'$ 的类似情形，使用“Bertini 定理”和 Zariski 的“\Term{connected}{连通}性定理”，实际上可以得到形式上明显更令人满意的结果；这曾促使我在口头讨论班中提出一个猜想，后来我已肯定地解决了它。因此在这里陈述如下。

\begin{SGAresult}{定理}{2.1}
设 $A$ 是\Term{complete}{完备}\Term{noetherian}{诺特}\Term{local-ring}{局部环}，$X$ 是其谱，$a$ 是 $X$ 的\Term{closed-point}{闭点}，$X'=X-\{a\}$。假设 $X$ 满足下列条件（其中 $k$ 为整数，$k\geq1$）：

$a_k)$ $X'$ 的\Term{irreducible-component}{不可约分支}的\Term{dimension}{维数}均 $\geq k+1$。
\SourcePage{148}

$b_k)$ $X'$ 在\Term{dimension}{维数} $\geq k$ 上\Term{connected}{连通}，即不能用一个\Term{dimension}{维数} $<k$ 的\Term{closed-subset}{闭子集}使 $X'$ 不\Term{connected}{连通}（参见 III 3.8）。

设 $m$ 为整数，$0\leq m\leq k$，且 $f_1,\ldots,f_m\in\mathfrak r(A)$，令
\[
B=A/\sum_i f_iA,\quad Y=\operatorname{Spec}(B)=V(f_1)\cap\cdots\cap V(f_m),
\quad Y'=X'\cap Y=Y-\{a\}.
\]
则 $Y$ 满足条件 $a_{k-m})$、$b_{k-m})$。特别地，对于 $\mathfrak r(A)$ 中任意 $m\leq k$ 个元素 $f_1,\ldots,f_m$ 的序列，$Y'=X'\cap V(f_1)\cap\cdots\cap V(f_m)$ \Term{connected}{连通}。
\end{SGAresult}

而且容易看出，若最后的结论成立（显然只须取 $m=k$），并排除 $X$ \Term{irreducible}{不可约}且\Term{dimension}{维数}为 $0$ 或 $1$ 的情形，则 $X'$ 的\Term{irreducible-component}{不可约分支}的\Term{dimension}{维数}均 $\geq k+1$，并且 $X'$ 在\Term{dimension}{维数} $\geq k$ 上\Term{connected}{连通}。因此，在某种意义下，2.1 是“最佳可能”的结果。

\begin{proof}
下面给出 2.1 的证明思路。只有 $Y$ 的条件 $b_{k-m})$ 需要处理。对于给定的 $k$，容易归结为 $X$ 整的情形，甚至（通过取\Term{normalization}{正规化}，它在 $X$ 上\Term{finite}{有限}）归结为 $X$ \Term{normal}{正规}的情形。若 $k=1$，从而 $\dim X'\geq2$，则 $X'$ 在其\Term{closed-point}{闭点}处\Term{depth}{深度} $\geq2$，可以应用本节开头回顾的结果，得到 $Y'=X'\cap V(f)$ \Term{connected}{连通}。

在 $k\geq1$ 的情形，假设定理已对 $k'<k$ 证明。对 $m$ 归纳，可归结为 $m=1$ 的情形，即验证对于 $f_1\in\mathfrak r(A)$，$X'\cap V(f_1)$ 在\Term{dimension}{维数} $\geq k-1$ 上\Term{connected}{连通}。如果不然，即它能被一个\Term{dimension}{维数} $<k-1$ 的 $Z'$ 断开，就存在序列 $f_2,\ldots,f_k$，使 $X'\cap V(f_1)\cap\cdots\cap V(f_k)$ 不\Term{connected}{连通}；而在这个序列中，$f_2\in\mathfrak r(A)$ 可以任意选择，只须在 $X'$ 的某个\Term{finite}{有限}\Term{subset}{子集} $F$（即 $Z'$ 的\Term{maximal-point}{极大点}\Term{set}{集合}）的任何点处都不为零。另一方面，利用 $X'$ \Term{normal}{正规}、从而满足 Serre 条件 $(S_2)$\footnote{参见 EGA IV 5.7.2。}，容易验证存在 $X'$ 的\Term{finite}{有限}\Term{subset}{子集} $F'$，使得 $f\in\mathfrak r(A)$ 及 $V(f)\cap F'=\varnothing$ 蕴含 $V(f)\cap X'$ 也满足 $(S_2)$。于是可选择 $f_2$ 使它在 $F$ 和 $F'$ 上都不为零，从而 $X'\cap V(f_2)$ 满足 $(S_2)$。但根据 Hartshorne 定理 III 3.6，$X'\cap V(f_2)$ 在\Term{codimension}{余维数} $1$ 上\Term{connected}{连通}，因而（由于 $X'\cap V(f_2)$ 的每个分支的\Term{dimension}{维数}均 $\geq k$）它在\Term{dimension}{维数} $\geq k-1$ 上\Term{connected}{连通}。把归纳假设用于 $V(f_2)=\operatorname{Spec}(A/f_2A)$，可得 $X'\cap V(f_2)\cap V(f_1)\cap V(f_3)\cap\cdots\cap V(f_k)$ \Term{connected}{连通}，而我们构造它时要求它不\Term{connected}{连通}，矛盾。
\end{proof}
指出几个有意思的推论。

\begin{SGAresult}{推论}{2.2}
设 $f:X\to Y$ 为\Term{local-scheme}{局部}\Term{noetherian}{诺特}\Term{prescheme}{预概形}的\Term{proper}{固有}\Term{morphism}{态射}，$Y$ 整，$y_0\in Y$，$y_1$ 为 $Y$ 的\Term{generic-point}{泛点}。假设：

 a) $Y$ 在 $y_0$ 处\Term{unibranch}{单分支}，且 $X$ 的每个\Term{irreducible-component}{不可约分支}都\Term{dominant}{支配} $Y$。
\SourcePage{149}

 b) $X_1$ 的\Term{irreducible-component}{不可约分支}的\Term{dimension}{维数}均 $\geq k+1$，且 $X_1$ 在\Term{dimension}{维数} $\geq k$ 上\Term{connected}{连通}。

则 $X_0$ 的\Term{irreducible-component}{不可约分支}的\Term{dimension}{维数}均 $\geq k+1$，且 $X_0$ 在\Term{dimension}{维数} $\geq k$ 上\Term{connected}{连通}。
\end{SGAresult}
\begin{proof}
事实上，Zariski \Term{connected}{连通}性定理（参见 EGA III 4.3.1）蕴含 $X_0$ \Term{connected}{连通}；为了证明一个\Term{dimension}{维数} $<k$ 的\Term{closed-subset}{闭子集}不能使它不\Term{connected}{连通}，可归结为证明：在满足 $\dim\overline{\{x\}}<k$ 的点 $x\in X_0$ 处，其\Term{local-ring}{局部环}的谱不能被 $x$ 断开。而这一点甚至无需假设 $f$ \Term{proper}{固有}、或 $Y$ 在 $y_0$ 处\Term{unibranch}{单分支}。为此，可归结为 $X$ 整且\Term{dominant}{支配} $Y$ 的情形；如果愿意，还可使 $Y$ 为 $\mathbf Z$ 上\Term{finite-type}{有限型}的\Term{affine-scheme}{仿射概形}，从而 $\mathcal O_{X,x}$ 在 $\mathcal O_{Y,y_0}$ 上满足\Term{dimension}{维数}公式的条件。利用此时\Term{normal}{正规}\Term{closure}{闭包}的\Term{finite}{有限}性，甚至可以假设 $X$ \Term{normal}{正规}；于是由 Nagata 的一个定理\footnote{参见 EGA IV 7.8.3 (i)、(ii)、(v)。}，$X'$ 的\Term{local-ring}{局部环} $\mathcal O_{X,x}$ 的\Term{completion}{完备化}仍\Term{normal}{正规}，因此（若 $\mathcal O_{X,x}$ 的\Term{dimension}{维数}为 $N$）$\operatorname{Spec}(\widehat{\mathcal O}_{X,x})$ 在\Term{dimension}{维数} $\geq N-1$ 上\Term{connected}{连通}。设 $n=\dim\mathcal O_{Y,y_0}$；考虑到 $\dim X_1\geq k+1$，$\operatorname{trdeg}k(x)/k(y)<k$ 蕴含
\[
\dim\mathcal O_{X,x}>n+(k+1)-k=n+1.
\]
取 $\mathcal O_{Y,y_0}$ 的一个\Term{system-of-parameters}{参数系} $f_1,\ldots,f_n$，把它们\Term{lifting}{提升}为 $\mathcal O_{X,x}$ 的元素，由 2.1 可见 $\operatorname{Spec}(\widehat{\mathcal O}_{X,x}/\sum f_i\widehat{\mathcal O}_{X,x})$ 在\Term{dimension}{维数} $\geq1$ 上\Term{connected}{连通}，即不能被其\Term{closed-point}{闭点}断开；\Term{equivalence}{等价}地，$\operatorname{Spec}(\widehat{\mathcal O}_{X_0,x})$ 不能被其\Term{closed-point}{闭点}断开；那么 $\operatorname{Spec}(\mathcal O_{X_0,x})$ 更是如此。
\end{proof}
如同通常的\Term{connected}{连通}性定理，2.2 还可以改为考虑\Term{geometric-fiber}{几何纤维}（即在\Term{residue-field}{剩余域}的\Term{algebraic-closure}{代数闭包}上取\Term{fiber}{纤维}），条件是 $Y$ 在 $y_0$ 处\Term{geometrically-unibranch}{几何单分支}，或者（除 $Y$ \Term{noetherian}{诺特}外不加其他假设）$f$ \Term{universally-open}{普遍开}。把这应用于 $Y$ 为域上\Term{projective-scheme}{射影概形} $\mathbf P^r_k$ 的\Term{duality}{对偶}\Term{scheme}{概形}的情形，就得到启发了 2.1 的那个整体结果的加强形式，即：

\begin{SGAresult}{推论}{2.3}
\footnote{编者注（6）：一个很漂亮的直接证明，见 Fulton W. \& Lazarsfeld R., “Connectivity and its applications in algebraic geometry”, 载 Algebraic geometry (Chicago, Ill., 1980), Lect. Notes in Math., 第 862 卷，Springer, Berlin--New York, 1981，26--92 页，定理 2.1。另参见原书第 155 页编者注（22）中引用的 [HL]。}
设 $X$ 是 $\mathbf P^r_k$ 的\Term{closed-subscheme}{闭子概形}（$k$ 为域），假设 $X$ 的\Term{irreducible-component}{不可约分支}的\Term{dimension}{维数}均 $\geq l+1$，且 $X$ 在\Term{dimension}{维数} $\geq l$ 上几何\Term{connected}{连通}。则对任意 $m$ 个 $\mathbf P^r_k$ 中\Term{hyperplane}{超平面}的序列 $H_1,\ldots,H_m$（$0\leq m\leq l-1$），$X\cap H_1\cap\cdots\cap H_m$ 满足把 $l$ 换成 $l-m$ 后的相同条件；特别地，它在\Term{dimension}{维数} $\geq l-1$ 上几何\Term{connected}{连通}。
\end{SGAresult}
对于给定一个不必为\Term{immersion}{浸入}的\Term{proper}{固有}\Term{morphism}{态射} $X\to\mathbf P^r_k$ 的情形，显然也可以修改这一命题；
\SourcePage{150}
2.1 也有类似推广（考虑 $X'$ 上的\Term{proper}{固有}\Term{scheme}{概形}）。事实上，利用通常的\Term{connected}{连通}性定理，把问题归结为\Term{finite-morphism}{有限态射}的情形后，这些命题即可形式地从这里给出的命题推出。

\begin{SGAresult}{推论}{2.4}
设 $A$ 是\Term{dimension}{维数} $\geq k+2$ 的\Term{complete}{完备}\Term{noetherian}{诺特}\Term{normal}{正规}\Term{local-ring}{局部环}。令 $X=\operatorname{Spec}(A)$、$X'=X-\{a\}$，设 $f_1,\ldots,f_k$ 是 $\mathfrak r(A)$ 中的元素，则 $Y'=X'\cap V(f_1)\cap\cdots\cap V(f_k)$ \Term{connected}{连通}，且 $\pi_1(Y')\to\pi_1(X')$ \Term{surjectivity}{满射}。
\end{SGAresult}
\begin{proof}
仿照 SGA 1 X 2.11 进行。
\end{proof}
以上只涉及\Term{connected}{连通}性问题。然而，在整体情形下，熟知的定理断言：对于\Term{irreducible}{不可约}\Term{projective-variety}{射影簇} $X\subset\mathbf P^r_k$，其中 $k$ \Term{algebraically-closed}{代数闭}，它与一个足够“一般”的\Term{hyperplane}{超平面} $H$ 的交\Term{irreducible}{不可约}（而不只是\Term{connected}{连通}）。这就是 Zariski 证明的 Bertini 定理，借助 Zariski \Term{connected}{连通}性定理，它又蕴含对于所有 $H$，$H\cap X$ \Term{connected}{连通}（虽不一定\Term{irreducible}{不可约}）。也可以反过来，先用 Lefschetz 型技巧证明后一结果，再由此推出 Bertini 定理：只须归结为 $X$ \Term{normal}{正规}的情形，并使用下述结果——对“足够一般”的 $H$，$X\cap H$ 也\Term{normal}{正规}。这提示了如下猜想。

\begin{SGAresult}{猜想}{2.5}
\footnote{编者注（7）：见下一条编者注。}
设 $A$ 是\Term{complete}{完备}\Term{noetherian}{诺特}\Term{normal}{正规}\Term{local-ring}{局部环}。证明存在非零的 $f\in\mathfrak r(A)$，使得 $Y'=X'\cap V(f)=Y-\{a\}$（其中 $Y=\operatorname{Spec}(A/fA)$）\Term{normal}{正规}（从而当 $\dim A\geq3$ 时由 2.1 得到它\Term{irreducible}{不可约}）。
\end{SGAresult}
为了做好这件事，还应当在某种适当的意义下证明，具有上述性质的元素 $f$ 甚至有“很多”，例如可以在\Term{maximal-ideal}{极大理想}的任意幂中选择 $f$。利用 Serre \Term{normal}{正规}性判据和前面对 Serre 性质 $(S_2)$ 的注记，可见若下述猜想有肯定解答，就会肯定地解答 2.5。

\begin{SGAresult}{猜想}{2.6}
\footnote{编者注（8）：现在文献中已有这个猜想的证明，因而如上所述，前一个猜想也应视为已获证明。此前还发表过 Flenner 和 Trivedi 的两次证明尝试，可惜均未成功。参见 Trivedi V., “Erratum: ‘A local Bertini theorem in mixed characteristic’”, Comm. Algebra 25 (1997), 第 5 期，1685--1686 页。不过，编者没有核实现在的证明是否已经完整。}
设 $A$ 是\Term{complete}{完备}\Term{noetherian}{诺特}\Term{local-ring}{局部环}，$U$ 是其谱 $X$ 的\Term{open-subset}{开子集}，$F$ 是 $X'=X-\{a\}$ 的\Term{finite}{有限}\Term{subset}{子集}。假设 $U$ \Term{regular}{正则}。证明存在 $f\in\mathfrak r(A)$，使 $V(f)\cap U$ \Term{regular}{正则}，且 $V(f)\cap F=\varnothing$。
\end{SGAresult}
有关一个“\Term{local-scheme}{局部} Bertini”型结果，见 Chow [2]。

\SourcePage{151}
\SGAsection{3}{与 $\pi_1$ 有关的问题}
这里同样有许多由整体结果或超越结果启发的问题。

\begin{SGAresult}{猜想}{3.1}
\footnote{编者注（9）：在所有\Term{finite-type}{有限型} $k$ \Term{scheme}{概形}都存在强\Term{resolution-of-singularities}{奇点消解}的假设下，类似命题对于\Term{separably-closed}{可分闭}域 $k$ 上的\Term{finite-type}{有限型}（\Term{connected}{连通}）\Term{scheme}{概形} $X$ 成立；特别地，当 $k$ \Term{characteristic}{特征}为零，或 $\dim X\leq2$ 时成立。为此，可用 Raynaud 夫人建立的 Lefschetz 型技巧归结为\Term{quasi-projective-morphism}{拟射影}曲面情形，参见前面的注释；见 SGA 7.1，定理 II.2.3.1。}
设 $A$ 是\Term{residue-field}{剩余域}\Term{algebraically-closed}{代数闭}的\Term{complete}{完备}\Term{noetherian}{诺特}\Term{local-ring}{局部环}，$X=\operatorname{Spec}(A)$，$X'=X-\{a\}$，其中 $a$ 为\Term{closed-point}{闭点}。假设 $X$ 的\Term{irreducible-component}{不可约分支}的\Term{dimension}{维数}均 $\geq2$，且 $X'$ \Term{connected}{连通}。

(i) 证明 $\pi_1(X')$ 在\Term{topology}{拓扑}意义下\Term{finitely-generated}{有限生成}。

(ii) 若 $p$ 是 $A$ 的\Term{residue-field}{剩余域} $k$ 的\Term{characteristic-exponent}{特征指数}，证明 $\pi_1(X')$ 的最大“阶与 $p$ 互素”的\Term{topology}{拓扑}\Term{quotient-group}{商群}是\Term{finitely-presented}{有限表示}的。
\end{SGAresult}
对于 (i)，使用\Term{descent}{下降}理论 SGA 1 IX 5.2 和定理 2.4，可以归结为 $A$ \Term{normal}{正规}且\Term{dimension}{维数}为 $2$ 的情形。在这种情形，Mumford [5] 在超越框架中开创的研究 $X'$ \Term{fundamental-group}{基本群}的系统方法，是对 $X$ \Term{resolution}{消解}奇点，即考虑一个\Term{projective}{射影}\Term{birational}{双有理}\Term{morphism}{态射} $Z\to X$，其中 $Z$ 整且\Term{regular}{正则}，并诱导\Term{isomorphism}{同构} $Z'=Z|_{X'}\to X'$；这样的 $Z$ 似乎总应存在，无论如何，在“\Term{equal-characteristic}{等特征}”情形，Abhyankar [1] 的方法确实证明了这一点。\footnote{现在 Abhyankar [8] 已在完全一般的情形下证明可以“\Term{resolution}{消解}”$A$。}设 $C$ 是 $Z\to X$ 在 $X$ 的\Term{closed-point}{闭点}上的\Term{fiber}{纤维}，它是\Term{residue-field}{剩余域} $k$ 上的一条\Term{algebraic-curve}{代数曲线}，由\Term{connected}{连通}性定理可知它\Term{connected}{连通}。于是 3.1 的解答看来与下列问题有关。

\begin{SGAresult}{问题}{3.2}
沿用上述记号，建立 $\pi_1(X')$ 与 $C$ 的\Term{topology}{拓扑}\Term{invariants}{不变量}，特别是其\Term{fundamental-group}{基本群}之间的关系（例如利用 SGA 1，定理 X 2.6，以显现出 $\pi_1(X')$ 的\Term{topology}{拓扑}\Term{finitely-generated}{有限生成}性）。
\end{SGAresult}
另一种方法是把 $A$ 看成\Term{dimension}{维数}为 $2$ 的\Term{complete}{完备}\Term{regular}{正则}\Term{local-ring}{局部环} $B$ 上的\Term{finite}{有限}\Term{algebra}{代数}，\Term{ramification}{分歧}沿着 $\operatorname{Spec}(B)=Y$ 中的一条曲线 $C$ 发生。这样就导向下列问题。

\begin{SGAresult}{问题}{3.3}
设 $A$ 是\Term{dimension}{维数}为 $2$、\Term{residue-field}{剩余域} $k$ \Term{algebraically-closed}{代数闭}的\Term{complete}{完备}\Term{regular}{正则}\Term{local-ring}{局部环}，$X$ 是其谱，$C$ 是 $X$ 的\Term{dimension}{维数}为 $1$ 的\Term{closed-subset}{闭子集}。定义\Term{immersion}{浸入}曲线 $C$ 的\Term{local-scheme}{局部}\Term{invariants}{不变量}，使其意义与\Term{residue-characteristic}{剩余特征}无关，并且知道这些\Term{invariants}{不变量}后，当 $k$ 的\Term{characteristic}{特征}为零时，能够用\Term{generator}{生成元}与关系计算
\SourcePage{152}
$X-C$ 的\Term{fundamental-group}{基本群}。证明，当 $k$ 的\Term{characteristic}{特征}为 $p>0$ 时，$X-C$ 的“\Term{tame}{驯}”\Term{fundamental-group}{基本群}是前一\Term{fundamental-group}{基本群}的商，并且这两个\Term{fundamental-group}{基本群}（\Term{characteristic}{特征}为 $0$ 与\Term{characteristic}{特征}为 $p>0$）有相同的阶与 $p$ 互素的最大\Term{quotient-group}{商群}。
\end{SGAresult}
当然，3.3 表明，在 3.1 中，也应当把 $X'$ 换成形如 $X-Y$ 的\Term{scheme}{概形}，其中 $Y$ 是 $X$ 的\Term{closed-subset}{闭子集}，并在包含它的 $X$ 的每个分支中\Term{codimension}{余维数}都 $\geq2$。若放弃对 $Y$ 的这一限制，仍应存在类似的\Term{finite}{有限}性结果，只须在 $Y$ 的那些\Term{codimension}{余维数}为 $1$ 的\Term{irreducible-component}{不可约分支}的\Term{maximal-point}{极大点}处，对\Term{ramification}{分歧}施加“\Term{tame}{驯}”一类的限制。

\begin{SGAresult}{问题}{3.4}
设 $A$ 是\Term{dimension}{维数}为 $2$、\Term{residue-field}{剩余域}\Term{algebraically-closed}{代数闭}的\Term{complete}{完备}\Term{noetherian}{诺特}\Term{local-ring}{局部环}。仍令 $X=\operatorname{Spec}(A)$，$X'=X-\{a\}$。在 $A$ 是\Term{complete-intersection}{完全交}的情形，寻找 $\pi_1(X')$ 特有的结构性质。
\end{SGAresult}
这个问题若有令人满意的解答，也许能解决下面这个老问题。

\begin{SGAresult}{猜想}{3.5}
\footnote{编者注（10）：截至 2004 年秋，这个问题仍未解决。}
在 $\mathbf P^3_k$（$k$ 为\Term{algebraically-closed}{代数闭}域）中找一条\Term{irreducible}{不可约}曲线，最好是\Term{nonsingular}{非奇异}的，使其不能在\Term{set}{集合}意义下表示为两个\Term{hypersurface}{超曲面}的交。
\end{SGAresult}
（Kneser [4] 证明，总可以把它表示为三个\Term{hypersurface}{超曲面}的交。）

\SGAsection{4}{与高阶 $\pi_i$ 有关的问题：复解析空间的局部和整体 Lefschetz 定理}
\footnote{编者注（11）：这些命题在评注（第 6 节）中得到更精确的表述。其中出现的猜想已经成为定理，见第 6 节的脚注。}
设 $X$ 是复数域 $\mathbf C$ 上\Term{locally-finite-type}{局部有限型}的\Term{scheme}{概形}，可以给它配以 $\mathbf C$ 上的\Term{analytic-space}{解析空间} $X^h$，从而得到\Term{homotopy}{同伦}与\Term{homology}{同调}\Term{invariants}{不变量} $\pi_i(X^h)$、$H_i(X^h)$、$H^i(X^h)$ 等。此外，已知 $X$ \Term{connected}{连通}当且仅当 $X^h$ \Term{connected}{连通}，因此有\Term{bijection}{双射}
\[
\pi_0(X^h)\longrightarrow\pi_0(X).
\]
同样，由于 $X$ 的每个\Term{etale-covering}{平展覆盖} $X'$ 都定义 $X^h$ 的\Term{etale-covering}{平展覆盖} $X'{}^h$，存在\Term{canonical}{典范}\Term{homomorphism}{同态}
\[
\pi_1(X^h)\longrightarrow\pi_1(X),
\]
\SourcePage{153}
利用 Grauert--Remmert 的一个定理，已知它把第二个群等同于第一个群对\Term{finite-index}{有限指数}\Term{subgroup}{子群}\Term{topology}{拓扑}的\Term{compactification}{紧化}（这只是说，$X'\mapsto X'{}^h$ 给出 $X$ 的\Term{etale-covering}{平展覆盖}\Term{category}{范畴}与 $X^h$ 的\Term{finite}{有限}\Term{etale-covering}{平展覆盖}\Term{category}{范畴}的\Term{equivalence}{等价}）。由此，本讨论班用纯\Term{algebra}{代数}方法得到的有关 $\pi_0(X)$ 和 $\pi_1(X)$ 的结果，蕴含有关 $\pi_0(X^h)$ 和 $\pi_1(X^h)$ 的超越性质的结果。此外，若 $X$ \Term{proper}{固有}，则熟知的\Term{exact-sequence}{正合列} $0\to\mathbf Z\to\mathbf C\to\mathbf C^*\to0$ 可用来证明 $X$ 的 Néron--Severi 群（即其 \Term{picard-group}{Picard 群}模去单位元的\Term{connected-component}{连通分支}所得的商）\Term{isomorphism}{同构}于 $H^2(X^h,\mathbf Z)$ 的一个\Term{subgroup}{子群}；在\Term{nonsingular}{非奇异} Kähler 情形，它就是记作 $H^{(1,1)}(X^h,\mathbf Z)$ 的\Term{subgroup}{子群}（$(1,1)$ 型的类）：
\[
\operatorname{Pic}(X)/\operatorname{Pic}^0(X)\subset H^2(X,\mathbf Z).
\]
于是，我们得到的有关 \Term{picard-group}{Picard 群}的信息，蕴含有关群 $H^2(X^h,\mathbf Z)$ 的信息，虽然确实很不完整。人们很想通过猜想来补足这些零碎结果。

Lefschetz 的一个经典定理 [7] 给出了与刚才所说一致的、十分精确的指引。它断言：若 $X$ 是\Term{dimension}{维数}为 $n$ 的\Term{nonsingular}{非奇异}\Term{irreducible}{不可约}\Term{projective}{射影}\Term{analytic-space}{解析空间}，$Y$ 是\Term{nonsingular}{非奇异}\Term{hyperplane-section}{超平面截面}，则嵌入
\[
Y^{n-1}\longrightarrow X^n
\]
诱导\Term{homomorphism}{同态}
\[
\pi_i(Y^{n-1})\longrightarrow\pi_i(X^n),
\]
它在 $i\leq n-2$ 时是\Term{isomorphism}{同构}，在 $i=n-1$ 时是\Term{epimorphism}{满态射}。由此，\Term{homology}{同调}（为明确起见，取整系数）上的\Term{homomorphism}{同态}
\[
H_i(Y^{n-1})\longrightarrow H_i(X^n)
\]
也满足类似命题；而在\Term{cohomology}{上同调}中，
\[
H^i(X^n)\longrightarrow H^i(Y^{n-1})
\]
在\Term{dimension}{维数} $i\leq n-2$ 时是\Term{isomorphism}{同构}，在\Term{dimension}{维数} $i=n-1$ 时是\Term{monomorphism}{单态射}。我们已在\Term{scheme}{概形}的框架内，对 $\pi_0$、$\pi_1$、$\operatorname{Pic}$ 得到这些结果的变体，而且这些变体在很大程度上不要求\Term{nonsingular}{非奇异}性，参见第 XII 讲。另外，在去掉\Term{nonsingular}{非奇异}性假设时，我们还实质地使用了这些整体 Lefschetz 定理的“\Term{local-scheme}{局部}”变体。所有这些提示下述问题，它们大概应当同时研究。\footnote{以下 4.1 至 4.3 的表述是暂定的。更令人满意的表述，见下文“第 XIII 讲评注”中的猜想 A 至 D，以及第 XIV 讲。}

\SourcePage{154}
\begin{SGAresult}{问题}{4.1}
设 $X$ 为\Term{analytic-space}{解析空间}，$Y$ 是 $X$ 的\Term{closed-analytic-subset}{闭解析子集}（或者仅仅是\Term{closed-subset}{闭子集}？）\footnote{编者注（12）：这个问题的含义不清楚；事实上，在这种情形，问题的命题本身似乎就没有意义，因为若不再假设 $Y$ 解析，就没有定义 $Y$ 在 $X$ 中的\Term{codimension}{余维数}。}，满足：对每个 $x\in Y$，\Term{local-ring}{局部环} $\mathcal O_{X,x}$ 都是\Term{complete-intersection}{完全交}。设 $n$ 为 $Y$ 在 $X$ 中的复\Term{codimension}{余维数}。\Term{canonical}{典范}\Term{homomorphism}{同态}
\[
\pi_i(X-Y)\longrightarrow\pi_i(X)
\]
是否在 $i\leq n-2$ 时为\Term{isomorphism}{同构}，在 $i=n-1$ 时为\Term{epimorphism}{满态射}？
\end{SGAresult}
在这个问题及下面的问题中，当然都默认选定了定义\Term{homotopy-group}{同伦群}所需的基点。为了陈述下一个问题，需要对\Term{analytic-space}{解析空间} $X$（更一般地，对\Term{local-scheme}{局部}道路\Term{connected}{连通}空间）和 $x\in X$ 定义\Term{local-scheme}{局部}\Term{invariants}{不变量} $\pi_i^x(X)$。\footnote{当 $i\geq2$ 时。对于 $i\leq1$ 的情形，参见下文第 6 节的评注，原书第 154 页。}为此，取区间 $[0,1]$ 到 $X$ 的非常值映射 $f$，使 $f(0)=x$，且 $t\ne0$ 时 $f(t)\ne x$（若 $x$ 不是\Term{isolated-point}{孤立点}，这样的映射存在）。于是，对 $x$ 的任意邻域 $U$，存在 $\varepsilon>0$，使 $0<t<\varepsilon$ 蕴含 $f(t)\in U$；而\Term{homotopy-group}{同伦群} $\pi_i(U-x,f(t))$ 本质上与 $t$ 无关（当 $t$ 变化时，它们由一个传递的\Term{isomorphism}{同构}系相联系），可记为 $\pi_i(U-x,f)$。于是令
\[
\pi_i^x(X)=\varprojlim_U\pi_{i-1}(U-x,f),
\]
\Term{inverse-limit}{逆极限}取遍 $x$ 的\Term{open-neighborhood}{开邻域}系 $U$。严格地说，这个\Term{limit}{极限}依赖于 $f$，应当记作 $\pi_i^x(X,f)$；不过可以验证，随 $f$ 变化，这些群彼此\Term{isomorphism}{同构}\footnote{至少当 $x$ 在其邻域内不使 $X$ 不\Term{connected}{连通}时如此，参见下文评注，原书第 154 页。}；更精确地，它们在从 $x$ 出发的上述类型的道路空间上构成一个\Term{local-system}{局部系统}。这些\Term{invariants}{不变量}是第 I 讲对 $X$ 上\Term{sheaf}{层} $F$ 引入的\Term{local-cohomology}{局部上同调}\Term{invariants}{不变量} $H^i_x(F)$ 的\Term{homotopy}{同伦}版本，应当起到 $X$ 模去 $X-x$ 的相对\Term{local-scheme}{局部}\Term{homotopy-group}{同伦群}的\Term{action}{作用}。设 $Y$ 是 $X$ 的一个\Term{topology}{拓扑}\Term{dimension}{维数} $\leq d$ 的\Term{closed-subset}{闭子集}；若对所有 $x\in Y$ 及 $i\leq n$ 这些\Term{invariants}{不变量}都消失，应当蕴含\Term{homomorphism}{同态}
\[
\pi_i(X-Y)\longrightarrow\pi_i(X)
\]
在 $i<n-d$ 时为\Term{bijection}{双射}，在 $i=n-d$ 时为\Term{surjectivity}{满射}。\footnote{修正后的表述，见下文评注，原书第 154 页。}从这个观点看，4.1（取 $Y$ 仅含一点）会蕴含一个纯\Term{local-scheme}{局部}性质的猜想，即当 $X$ 在 $x$ 处为\Term{dimension}{维数}为 $n$ 的\Term{complete-intersection}{完全交}时，
\[
\pi_i^x(X)=0\qquad\text{当 }i\leq n-1.
\]
\SourcePage{155}
作为\Term{local-scheme}{局部}\Term{invariants}{不变量} $\pi_i^x(X)$ 的例子，注意：若 $x$ 是 $X$ 的一个复\Term{dimension}{维数}为 $n$ 的\Term{nonsingular}{非奇异}点，则
\[
\pi_i^x(X)=\pi_{i-1}(S^{2n-1}),
\]
其中 $S^{2n-1}$ 表示\Term{dimension}{维数}为 $2n-1$ 的球面。特别地，在此情形 $i\leq2n-1$ 时 $\pi_i^x(X)=0$；这对应于下述事实：若从\Term{topological-manifold}{拓扑流形} $X$ 中去掉一个\Term{codimension}{余维数} $\geq m$ 的\Term{closed-subset}{闭子集} $Y$，则 $\pi_i(X-Y)\to\pi_i(X)$ 在 $i\leq m-2$ 时是\Term{isomorphism}{同构}，在 $i=m-1$ 时是\Term{epimorphism}{满态射}。

有了这些，提出：
\begin{SGAresult}{问题}{4.2}
设 $X$ 是\Term{analytic-space}{解析空间}，$x\in X$，$t$ 是 $\mathcal O_X$ 的一个在 $x$ 处为零的\Term{section}{截面}，$Y$ 是 $t$ 的零点集。假设满足下述条件：

 a) $t$ 在 $x$ 处\Term{regular}{正则}（即在 $x$ 处不是\Term{zero-divisor}{零因子}；不过这个假设也许多余）。

 b) 在 $X-Y$ 中接近 $x$ 的点 $x'$ 处，$\mathcal O_{X,x'}$ 是\Term{complete-intersection}{完全交}（若 4.1 为真，这个假设应当能替换为如下更一般的假设：对于上述 $x'$，当 $i\leq n-1$ 时，$\pi_i^{x'}(X)=0$）。

 c) 在 $Y-\{x\}$ 中接近 $x$ 的点 $y$ 处，有
\[
\operatorname{prof}\mathcal O_{X,y}\geq n
\]
（例如，只须 $\operatorname{prof}\mathcal O_{X,x}\geq n$ 即可）。

在这些条件下，\Term{canonical}{典范}\Term{homomorphism}{同态}
\[
\pi_i^x(Y)\longrightarrow\pi_i^x(X)
\]
是否在 $i\leq n-2$ 时为\Term{isomorphism}{同构}，在 $i=n-1$ 时为\Term{epimorphism}{满态射}？
\end{SGAresult}
最后，下面是 4.2 的一个整体变体，它应当可以通过考察\Term{projecting-cone}{投影锥}的原点从 4.2 推出，并且将推广经典 Lefschetz 定理。

\begin{SGAresult}{问题}{4.3}
设 $X$ 是\Term{projective}{射影}\Term{analytic-space}{解析空间}，配有\Term{ample}{丰沛}\Term{invertible-module}{可逆模} $L$；$t$ 是 $L$ 的\Term{section}{截面}，$Y$ 是 $t$ 的零点集。假设：

 a) $t$ 是\Term{regular}{正则}\Term{section}{截面}（这个假设也许多余）。

 b) 对每个 $x\in X-Y$，$\mathcal O_{X,x}$ 是\Term{complete-intersection}{完全交}（应当可以用 $i\leq n-1$ 时 $\pi_i^x(X)=0$ 来代替）。

 c) 对每个 $x\in Y$，$\operatorname{prof}\mathcal O_{X,x}\geq n$。

在这些条件下，\Term{homomorphism}{同态}
\[
\pi_i(Y)\longrightarrow\pi_i(X)
\]
是否在 $i\leq n-2$ 时为\Term{isomorphism}{同构}，在 $i=n-1$ 时为\Term{epimorphism}{满态射}？
\end{SGAresult}

\SourcePage{156}
我们把陈述类似的\Term{cohomology}{上同调}性质的猜想留给读者\footnote{相应的\Term{scheme}{概形}理论中的结果，参见第 XIV 讲。}；届时，假设与结论都涉及\Term{local-cohomology}{局部上同调}\Term{invariants}{不变量}（系数取为某个给定群）。无论从\Term{homology}{同调}观点还是从\Term{homotopy}{同伦}观点来看，关键结果似乎都应当是 4.2，其中假设 b) 取括号中给出的较一般形式。

我们在超越框架中陈述这些猜想，希望能引起\Term{topology}{拓扑}学家的兴趣，使他们相信“Lefschetz”型问题远未完结。当然，凭借 M. Artin 最近的工作，我们眼下就将拥有良好的\Term{scheme}{概形}\Term{cohomology}{上同调}理论（取\Term{finite}{有限}系数）；于是，同样的问题也出现在\Term{scheme}{概形}的框架中，而且很难怀疑，它们将在不久的将来得到肯定回答。\footnote{见上一注。}

\SGAsection{5}{与局部 Picard 群有关的问题}
第一个基本问题，由 Mumford [5] 首先在一个特殊情形中指出，表述如下。设 $A$ 是\Term{residue-field}{剩余域}为 $k$ 的\Term{complete}{完备}\Term{local-ring}{局部环}，$X=\operatorname{Spec}(A)$，$U=\operatorname{Spec}(A)-\{a\}$，其中 $a$ 是 $A$ 的\Term{maximal-ideal}{极大理想}，即 $\operatorname{Spec}(A)$ 的\Term{closed-point}{闭点}。我们希望构造一个由 $k$ 上\Term{locally-algebraic-group}{局部代数群} $G_i$ 构成的严格\Term{projective-system}{逆系} $G$，以及自然\Term{isomorphism}{同构}
\[
\operatorname{Pic}(U)\simeq G(k),\tag{+}
\]
其中当然令 $G(k)=\varprojlim G_i(k)$。启发式地说，就是要在群 $\operatorname{Pic}(U)$ 上“赋予\Term{algebraic-group}{代数群}结构”（或者至少在适当意义下赋予\Term{algebraic-group}{代数群}的\Term{inverse-limit}{逆极限}结构）。

显然，这样提出的问题还不够精确，因为给定\Term{isomorphism}{同构} (+) 远远不能刻画\Term{pro-object}{逆对象} $G$。若 $A$ 包含一个仍记为 $k$ 的\Term{coefficient-field}{代表域}，就可以把问题明确为：要求对变动的 $k$ 的\Term{group-extension}{扩张} $k'$，有关于 $k'$ \Term{functorial}{函子性}的\Term{isomorphism}{同构}
\[
\operatorname{Pic}(U')\simeq G(k'),\tag{+'}
\]
其中 $U'$ 是 $\operatorname{Spec}(A')$ 中与 $U$ 类似的\Term{open-set}{开集}，$A'=A\widehat\otimes_k k'$。即使 $A$ 没有\Term{coefficient-field}{代表域}，只要 $k$ 完全，也可以类似地进行，因为此时能够“通过\Term{residue-field}{剩余域}\Term{group-extension}{扩张} $k'/k$”\Term{functorial}{函子性}地构造 $A'$。另外，当 $A$ 有\Term{coefficient-field}{代表域}时，在 $\operatorname{Pic}(U)$ 上得到的\Term{algebraic-structure}{代数结构}将实质地依赖于这个\Term{coefficient-field}{代表域}的选择（从\Term{elliptic-curve}{椭圆曲线}的\Term{projecting-cone}{投影锥}就已能看出），因此，似乎必须从 Greenberg [3] 意义下的“$k$ 上的\Term{algebra}{代数}环的\Term{inverse-limit}{逆极限}”出发，
\SourcePage{157}
才能定义\Term{pro-object}{逆对象} $G$。而在没有给定\Term{coefficient-field}{代表域}的情形，也可以设想最终只得到一个由 Serre 意义下的拟\Term{algebraic-group}{代数群}，或者更确切地说，拟\Term{locally-algebraic-group}{局部代数群}构成的\Term{projective-system}{逆系}（得到的群 $G_i$ 一般并非 $k$ 上\Term{finite-type}{有限型}，而只在 $k$ 上\Term{locally-finite-type}{局部有限型}）。甚至可能，一般在 $\operatorname{Pic}(U)$ 上只能得到更弱的结构，类似于 Néron [6] 在其\Term{local-scheme}{局部}域上定义的\Term{abelian-variety}{阿贝尔簇}\Term{degeneration}{退化}理论中遇到的那些结构。

Mumford 也提出了研究这个问题的一种方法：对 $X$ \Term{resolution}{消解}奇点，即考虑\Term{projective}{射影}\Term{birational}{双有理}\Term{morphism}{态射} $Y\to X$，其中 $Y$ \Term{regular}{正则}。当 $U$ \Term{regular}{正则}（即 $a$ 是\Term{isolated-singularity}{孤立奇点}）时，常可找到 $Y$，使 $Y|_U=V\to U$ 为\Term{isomorphism}{同构}。在这种情形，有
\[
\operatorname{Pic}(U)\simeq\operatorname{Pic}(V)
 \simeq\operatorname{Pic}(Y)/\operatorname{Im}\mathbf Z^I,
\]
其中 $I$ 是\Term{fiber}{纤维} $Y_a$ 的\Term{irreducible-component}{不可约分支}\Term{set}{集合}（每个分支都定义 $\operatorname{Pic}(Y)$ 的一个元素，因为 $Y$ \Term{regular}{正则}，所以它是\Term{locally-principal-divisor}{局部主除子}）。另一方面，利用\Term{formal-geometry}{形式几何}的技巧 EGA III 4 和 5，特别是\Term{existence-theorem}{存在定理}，得到
\[
\operatorname{Pic}(Y)\simeq\varprojlim\operatorname{Pic}(Y_n),
\]
其中 $Y_n=Y\otimes_A A_n$，$A_n=A/\mathfrak m^{n+1}$。当 $A$ 有\Term{coefficient-field}{代表域} $k$ 时，可以使用 $k$ 上\Term{projective-scheme}{射影概形} $Y_n$ 的 \Term{picard-scheme}{Picard 概形}理论，因而有
\[
\operatorname{Pic}(Y_n)\simeq\operatorname{Pic}_{Y_n/k}(k).
\]
这就给出了\Term{locally-algebraic-group}{局部代数群} $\operatorname{Pic}_{Y_n/k}/\operatorname{Im}\mathbf Z^I$ 构成的\Term{projective-system}{逆系}，它正是所求的\Term{projective-system}{逆系}。\footnote{编者注（13）：Boutot 的结果极大地澄清了这个问题，见 Boutot J.-F., Schéma de Picard local, Lect. Notes in Math., 第 632 卷，Springer, Berlin, 1978。特别地，若 $A$ 是\Term{depth}{深度} $\geq2$ 的\Term{complete}{完备}（\Term{noetherian}{诺特}）\Term{local-scheme}{局部} $k$ \Term{algebra}{代数}，且 $H^2_{\mathfrak m}(A)$ 在 $k$ 上\Term{finite}{有限}维，则\Term{local-picard-group}{局部 Picard 群}是 $k$ 上\Term{locally-finite-type}{局部有限型}的\Term{group-scheme}{群概形}，原点处的\Term{tangent-space}{切空间}为 $H^2_{\mathfrak m}(A)$。若 $A$ 还\Term{normal}{正规}且\Term{dimension}{维数} $\geq3$，则 Serre \Term{normal}{正规}性判据 XI 3.11 与推论 V 3.6 保证所需的\Term{finite}{有限}性，因而保证\Term{local-picard-scheme}{局部 Picard 概形}的存在。另参见 Lipman J., “The Picard group of a scheme over an Artin ring”, Publ. Math. Inst. Hautes Études Sci. 46 (1976), 15--86 页，其方法更接近上面勾勒的 Grothendieck 的方法。}在这里考虑的情形，还能看出（利用 $a$ 是\Term{isolated-singularity}{孤立奇点}），这个\Term{projective-system}{逆系}的\Term{universal}{普遍}像\Term{subgroup}{子群}的\Term{connected-component}{连通分支}构成一个\Term{essentially-constant-ind-object}{本质常值}的\Term{projective-system}{逆系}，因此这里得到了一个\Term{locally-algebraic-group}{局部代数群} $G$ 作为问题的解答。若还假设 $A$ \Term{normal}{正规}且\Term{dimension}{维数}为 $2$，则 Mumford 的一条注记（即 $Y_a$ 的各分支在 $X$ 中的\Term{intersection-matrix}{交矩阵}负定\footnote{编者注（14）：Mumford D., “The topology of normal singularities of an algebraic surface and a criterion for simplicity”, Publ. Math. Inst. Hautes Études Sci. 9 (1961), 5--22 页。}）蕴含 $G$ 甚至是
\SourcePage{158}
\Term{algebraic-group}{代数群}，即 $k$ 上\Term{finite-type}{有限型}（其\Term{connected-component}{连通分支}数等于刚才所考虑的\Term{intersection-matrix}{交矩阵}的行列式）。

反之，若 $a$ 不是\Term{isolated-singularity}{孤立奇点}，则从例子（$A$ \Term{dimension}{维数}为 $2$）可以看出，得到的是\Term{algebraic-group}{代数群}的\Term{projective-system}{逆系}，不能化为单个\Term{algebraic-group}{代数群}。

一旦有了良好的“\Term{local-picard-scheme}{局部 Picard 概形}”概念，就应当加强\Term{parafactorial}{副阶乘}性的概念：若不仅 $A$ 乃至 $\widehat A$ 是\Term{parafactorial}{副阶乘}的，而且\Term{local-picard-scheme}{局部 Picard 概形} $G(\widehat A)$ 是\Term{trivial-bundle}{平凡}群，就称 $A$ 为“\Term{geometrically-parafactorial}{几何副阶乘}的”（当\Term{residue-field}{剩余域}不是\Term{algebraically-closed}{代数闭}域时，这比说 $G$ 的 $k$ \Term{rational-point}{有理点}只有单位元要强）。回想存在\Term{dimension}{维数}为 $2$ 的\Term{complete}{完备}\Term{normal}{正规}\Term{local-ring}{局部环}，它们是\Term{factorial}{阶乘}环，却有不是\Term{factorial}{阶乘}环的\Term{finite}{有限}\Term{etale}{平展}\Term{algebra}{代数}\footnote{编者注（15）：在研究\Term{vector-bundle}{向量丛}的模空间时，自然会出现\Term{henselization}{亨泽尔化}不为\Term{factorial}{阶乘}环的\Term{factorial}{阶乘}环。例如参见 Drézet J.-M., “Groupe de Picard des variétés de modules de faisceaux semi-stables sur $\mathbf P^2$”, 载 Singularities, representation of algebras, and vector bundles (Lambrecht, 1985), Lect. Notes in Math., 第 1273 卷，Springer, Berlin, 1987，337--362 页。严格说来，Drézet 证明的是\Term{completion}{完备化}不为\Term{factorial}{阶乘}环；但实际上其证明也给出\Term{henselization}{亨泽尔化}的结果：关键在于 Luna 的\Term{etale}{平展}切片定理（Luna D., “Slices étales”, 载 Sur les groupes algébriques, Mém. Soc. math. France, 第 33 卷，Société mathématique de France, Paris, 1973，81--105 页），它在半\Term{stable}{稳定}点附近，对\Term{etale-topology}{平展拓扑}\Term{local-scheme}{局部}地描述\Term{invariants}{不变量}几何意义下的商的\Term{local-ring}{局部环}。}，就能意识到加强\Term{parafactorial}{副阶乘}性概念的必要性。于是，一个“\Term{geometrically-factorial}{几何阶乘}的”\Term{local-ring}{局部环}应是\Term{normal}{正规}环 $A$，使其所有\Term{dimension}{维数} $\geq2$ 的\Term{localization}{局部化}都\Term{geometrically-parafactorial}{几何副阶乘}；或者更好地说，要求 $\widehat A$ 的各个\Term{localization}{局部化}都是\Term{parafactorial}{副阶乘}的。\footnote{关于“\Term{geometrically-factorial}{几何阶乘}的”\Term{local-ring}{局部环}更灵活的概念，见原书第 152 页评注。}当然，为这些概念找到不依赖于尚待建立的\Term{local-picard-scheme}{局部 Picard 概形}理论的“好”定义，会是有意思的。\footnote{编者注（16）：见原书第 152 页：若一个\Term{local-ring}{局部环}的\Term{strict-henselization}{严格亨泽尔化}为\Term{factorial}{阶乘}的（分别地，\Term{parafactorial}{副阶乘}的），就称此\Term{local-ring}{局部环}为\Term{geometrically-factorial}{几何阶乘}的（分别地，\Term{geometrically-parafactorial}{几何副阶乘}的）。}

无论如何，若要得到下述类型的命题，似乎很可能需要这些概念。设 $A$ 是“好环”（例如 $\mathbf Z$ 上或\Term{complete}{完备}\Term{local-ring}{局部环}上的\Term{finite-type}{有限型}\Term{algebra}{代数}，例如域上的\Term{finite-type}{有限型}\Term{algebra}{代数}）。设 $U$ 是所有满足 $\mathcal O_{X,x}$ “\Term{geometrically-factorial}{几何阶乘}”的点 $x\in X=\operatorname{Spec}(A)$ 的\Term{set}{集合}，那么 $U$ 是否为\Term{open-set}{开集}？又或者：设 $f:X\to Y$ 为\Term{finite-type}{有限型}\Term{flat}{平坦}\Term{morphism}{态射}，其中 $Y$ \Term{local-scheme}{局部}\Term{noetherian}{诺特}；设 $U$ 是所有满足 $\mathcal O_{X_{f(x)},x}$ “\Term{geometrically-factorial}{几何阶乘}”的点 $x\in X$ 的\Term{set}{集合}，那么至少在 $f$ 满足一些良好的附加条件时，$U$ 是否为\Term{open-set}{开集}？我怀疑，使用通常的\Term{factorial}{阶乘}环概念，能否有这种类型的正确命题。

\SourcePage{159}
我们在这里通过一个特殊情形，提出了研究“变动的”\Term{local-ring}{局部环}的几何性质的问题，例如当 $x$ 遍历\Term{prescheme}{预概形} $X$ 时，研究 $\mathcal O_{X,x}$。例如，当 $X$ 是域上\Term{finite-type}{有限型}\Term{scheme}{概形}时，已知\footnote{编者注（17）：见 EGA IV.16。}$X$ 上存在\Term{finite}{有限}\Term{algebra}{代数}构成的\Term{projective-system}{逆系} $\mathcal P^n_{X/k}$（通过沿\Term{diagonal}{对角线}\Term{completion}{完备化} $X\times_kX$ 得到），它在每个 $k$ \Term{rational-point}{有理点} $x\in X$ 处的\Term{fiber}{纤维}\Term{isomorphism}{同构}于 $\mathcal O_{X,x}/\mathfrak m_X^{n+1}$ 的\Term{projective-system}{逆系}。因此，很自然地把变动的 $x$ 处\Term{local-ring}{局部环} $\mathcal O_{X,x}$ 的\Term{completion}{完备化}的研究，与 $\mathcal P^n$ 给出的“\Term{complete}{完备}\Term{local-ring}{局部环}的\Term{algebra}{代数}族”的研究联系起来；注意，对任意 $x\in X$（不论是否为 $k$ \Term{rational-point}{有理点}），都能得到\Term{complete}{完备}\Term{local-ring}{局部环}
\[
\mathcal P^\infty(x)=\varprojlim\mathcal P^n(x)
\]
（其中 $\mathcal P^n(x)$ 为\Term{reduced}{约化}\Term{fiber}{纤维} $\mathcal P^n\otimes_{\mathcal O_{X,x}}k(x)$）。例如，这样附着于\Term{generic-point}{泛点}的\Term{complete}{完备}环会特别值得关注；可以期待，它的\Term{algebra}{代数}几何性质（例如用其\Term{local-picard-group}{局部 Picard 群}、\Term{local-scheme}{局部}\Term{homotopy-group}{同伦群}或\Term{local-scheme}{局部}\Term{homology}{同调}群表述的性质），本质上就是 $x$ 遍历某个合适的\Term{dense-open-set}{稠密开集} $U$ 时，\Term{completion}{完备化} $\widehat{\mathcal O}_{X,x}$ 的性质。

一般地，可以提出同时研究由给定\Term{scheme}{概形} $X$ 上\Term{finite}{有限}\Term{algebra}{代数}构成的进制\Term{projective-system}{逆系} $(\mathcal P^n)$ 所得到的这些\Term{complete}{完备}\Term{local-ring}{局部环}。在一些\Term{regularity}{正则性}条件（如 $\mathcal P^n$ 的\Term{flat}{平坦}性）下，似乎很可能发现\Term{local-scheme}{局部}\Term{homotopy-group}{同伦群}来自 $X$ 上\Term{finite}{有限}\Term{group-scheme}{群概形}的一个\Term{projective-system}{逆系}，而\Term{local-picard-group}{局部 Picard 群}也会有类似结果。对于后者，第一个值得研究的有意思的情形，是从一张具有\Term{singular}{奇异}曲线的\Term{algebra}{代数}曲面 $X$ 出发，用定义在\Term{singular}{奇异}轨迹上的合适的\Term{group-scheme}{群概形}的\Term{inverse-limit}{逆极限}，研究沿这些曲线变动的点处的\Term{local-picard-scheme}{局部 Picard 概形}。

\SGAsection{6}{评注}
\footnote{写于 1963 年 3 月。}
\Term{scheme}{概形}的“\Term{etale}{平展}\Term{cohomology}{上同调}”观点及这一理论最近的进展，使我们既能明确又能推广前面提出的某些问题。有关“\Term{topology}{拓扑}”和“\Term{scheme}{概形}的\Term{etale-topology}{平展拓扑}”的概念，读者可参阅 M. Artin, Grothendieck Topologies, Harvard University, 1962（油印笔记）。\footnote{或者最好参阅 SGA 4。}

这一理论通过比传统的“Zariski \Term{topology}{拓扑}”所提供的更精细的\Term{localization}{局部化}概念，使人们特别关注\Term{strictly-local-ring}{严格局部环}，即\Term{residue-field}{剩余域}\Term{separably-closed}{可分闭}的\Term{henselian}{亨泽尔}\Term{local-ring}{局部环}。
\SourcePage{160}
对于\Term{residue-field}{剩余域}为 $k$ 的任意\Term{local-ring}{局部环} $A$，以及 $k$ 的任意\Term{separable-closure}{可分闭包} $k'$，可以找到 $A$ 到一个\Term{strictly-local-ring}{严格局部环} $A'$ 的\Term{local-scheme}{局部}\Term{homomorphism}{同态}，$A'$ 称为 $A$ 的\Term{strict-local-closure}{严格局部闭包}，\Term{residue-field}{剩余域}为 $k'$，并满足一个显然的\Term{universal-property}{泛性质}。$A'$ 是\Term{henselian}{亨泽尔}的，在 $A$ 上\Term{flat}{平坦}，且 $A'\otimes_Ak\simeq k'$；它为\Term{noetherian}{诺特}环当且仅当 $A$ 为\Term{noetherian}{诺特}环。（参见所引文献，第 III 章第 4 节。）\footnote{或者 EGA IV 18.8。}若 $X$ 为\Term{prescheme}{预概形}，$x$ 为 $X$ 的点，$x'$ 为 $x$ 上方的一个点，是 $k=k(x)$ 的一个\Term{separable-closure}{可分闭包} $k'$ 的谱，则自然把 $X$ 在 $x'$ 处的\Term{strictly-local-ring}{严格局部环} $\mathcal O'_{X,x'}$ 定义为通常\Term{local-ring}{局部环} $\mathcal O_{X,x}$ 关于\Term{residue-field}{剩余域}\Term{group-extension}{扩张} $k'/k$ 的\Term{strict-local-closure}{严格局部闭包}。从\Term{etale-topology}{平展拓扑}观点看，正是 $X$ 的“几何”点处的\Term{strictly-local-ring}{严格局部环}，才应当反映\Term{prescheme}{预概形} $X$ 的\Term{local-scheme}{局部}性质。在很多方面，它们还起到原先由 $X$ 的\Term{local-ring}{局部环}的\Term{completion}{完备化}（比方说在\Term{residue-field}{剩余域}\Term{algebraically-closed}{代数闭}的点处）所起的\Term{action}{作用}，同时又与 $X$ “更接近”，使得向“邻近点”的过渡更为容易。

因此，应当对\Term{noetherian}{诺特}\Term{henselian}{亨泽尔}\Term{local-ring}{局部环}（分别地，对\Term{noetherian}{诺特}\Term{strictly-local-ring}{严格局部环}），重新提出通常针对\Term{complete}{完备}\Term{local-ring}{局部环}（有时还限于\Term{residue-field}{剩余域}\Term{algebraically-closed}{代数闭}的情形）提出的许多问题。例如，第 2、3 节提出的\Term{topology}{拓扑}问题，更一般地也适用于\Term{strictly-local-ring}{严格局部环}。另外，可以猜想，对于“好”的\Term{strictly-local-ring}{严格局部环}，\Term{canonical}{典范}\Term{morphism}{态射} $\operatorname{Spec}(\widehat A)\to\operatorname{Spec}(A)$ 的\Term{geometric-fiber}{几何纤维}具有某些\Term{simply-connected}{单连通}性和\Term{acyclic}{无环}性性质；这些性质会表明，对于许多“\Term{topology}{拓扑}”性质，在环 $A$ 上证明，或在其\Term{completion}{完备化} $\widehat A$ 上证明，是一回事。这个方向上已有的某些结果\footnote{参见 M. Artin 在 SGA 4 XIX 中的工作。}，使人有望很快在这一方向获得完整的结果。

\Term{etale}{平展}\Term{localization}{局部化}的概念，给出了\Term{local-ring}{局部环}“\Term{geometrically-parafactorial}{几何副阶乘}”或“\Term{geometrically-factorial}{几何阶乘}”概念的一个看来合理的定义（第 5 节原书第 150 页已经指出需要这种概念）：一个\Term{local-ring}{局部环}若其\Term{strict-local-closure}{严格局部闭包}\Term{parafactorial}{副阶乘}，分别地，\Term{factorial}{阶乘}，就这样称呼它。这类假设确实自然地出现在\Term{prescheme}{预概形}的\Term{etale}{平展}\Term{cohomology}{上同调}研究中。\footnote{编者注（18）：例如见 Strano R., “The Brauer group of a scheme”, Ann. Mat. Pura Appl. (4) 121 (1979), 157--169 页；其中\Term{scheme}{概形} $X$ 的\Term{local-ring}{局部环}的\Term{geometrically-parafactorial}{几何副阶乘}性假设，有时可以用来证明 $X$ 的 \Term{brauer-group}{Brauer 群}（用 Azumaya \Term{algebra}{代数}计算）与 $X$ 的\Term{cohomology}{上同调} \Term{brauer-group}{Brauer 群}一致。}例如，若 $X$ 是\Term{local-scheme}{局部}\Term{noetherian}{诺特}\Term{prescheme}{预概形}，其\Term{strictly-local-ring}{严格局部环}均为\Term{factorial}{阶乘}环（即通常的\Term{local-ring}{局部环}都“\Term{geometrically-factorial}{几何阶乘}”），则可以证明，$i\geq2$ 时 $H^i(X_{\text{ét}},\mathbf G_m)$
\SourcePage{161}
都是\Term{torsion-group}{挠群}（有时由此可以用以 $n$ 次\Term{roots-of-unity}{单位根群} $\boldsymbol\mu_n$ 为系数的\Term{cohomology}{上同调}群来表示这些群）；若 $X$ 整，\Term{fraction-field}{分式域}为 $K$，则自然\Term{homomorphism}{同态}
\[
H^2(X_{\text{ét}},\mathbf G_m)\longrightarrow H^2(K,\mathbf G_m)=\operatorname{Br}(K)
\]
是\Term{injectivity}{单射}。\footnote{编者注（19）：\Term{brauer-group}{Brauer 群}与 \Term{picard-group}{Picard 群}的关系十分密切。在这里引用 Saito 关于曲面的下述结果，第一个是\Term{local-scheme}{局部}的，另一个是整体的，见 Saito S., “Arithmetic on two-dimensional local rings”, Invent. Math. 85 (1986), 第 2 期，379--414 页。设 $A$ 是二维\Term{excellent}{优秀}\Term{local-ring}{局部环}，\Term{normal}{正规}且\Term{henselian}{亨泽尔}，\Term{residue-field}{剩余域}\Term{finite}{有限}；$X$ 是 $\operatorname{Spec}(A)$ 中\Term{closed-point}{闭点}的补集。则有\Term{torsion-group}{挠群}的\Term{perfect-duality}{完美对偶} $\operatorname{Pic}(X)\times\operatorname{Br}(X)\to\mathbf Q/\mathbf Z$；这里 $X$ 的 \Term{brauer-group}{Brauer 群}指\Term{cohomology}{上同调} \Term{brauer-group}{Brauer 群} $\operatorname{Br}(X)=H^2_{\text{ét}}(X,\mathbf G_m)$。在整体情形，Lichtenbaum 的一个结果有如下推广，参见 Lichtenbaum S., “Duality theorems for curves over $p$-adic fields”, Invent. Math. 7 (1969), 120--136 页：设 $k$ 为\Term{residue-field}{剩余域}\Term{finite}{有限}的\Term{complete}{完备}\Term{discrete-valuation-ring}{离散赋值环} $\mathcal O$ 的\Term{fraction-field}{分式域}，$X$ 为 $k$ 上\Term{projective}{射影}、\Term{smooth}{光滑}且几何\Term{complete}{完备}的曲线。给群 $\operatorname{Pic}^0(X)$ 赋予从 $k$ 的\Term{adic-topology}{进拓扑}诱导的\Term{topology}{拓扑}，并令 $\operatorname{Pic}(X)$ 为使 $\operatorname{Pic}^0(X)$ 成为开\Term{subgroup}{子群}的\Term{topological-group}{拓扑群}。则有\Term{topological-group}{拓扑群}的\Term{perfect-duality}{完美对偶} $\operatorname{Pic}(X)\times\operatorname{Br}(X)\to\mathbf Q/\mathbf Z$。注意，这个关于曲线的命题，当然是通过考虑 $X$ 在 $\mathcal O$ 上的\Term{regular}{正则}（\Term{proper}{固有}而\Term{flat}{平坦}的）模型来证明的：它是关于曲面的结果。}有例子表明，若只假设 $X$ \Term{factorial}{阶乘}而非\Term{geometrically-factorial}{几何阶乘}，哪怕 $X$ 是\Term{local-scheme}{局部}的，这些结论也可能不成立。\footnote{参见 A. Grothendieck, Le groupe de Brauer II（Séminaire Bourbaki 第 297 讲，1965 年 11 月），特别是 1.8 和 1.11 b。}

关于 3.4 中提出的\Term{local-scheme}{局部}与整体 Lefschetz 型问题，以及它们在\Term{scheme}{概形}理论中的类似问题，这些问题的\Term{homology}{同调}版本已经清楚得多，一切都形式地来自三个一般定理。第一个有关某些\Term{affine-scheme}{仿射概形}（分别地，Stein 空间）的\Term{cohomological-dimension}{上同调维数}，例如\Term{algebraically-closed}{代数闭}域上\Term{finite-type}{有限型}\Term{affine-scheme}{仿射概形} $X$：它们的\Term{cohomological-dimension}{上同调维数} $\leq\dim X$（“\Term{affine}{仿射} Lefschetz 定理”）\footnote{参见 SGA 4 XIV。}；第二个是\Term{projective-morphism}{射影态射}的\Term{cohomology}{上同调}（\Term{discrete}{离散}系数）的\Term{duality}{对偶}定理\footnote{参见 SGA 4 XVIII。}；最后一个则是类似性质的\Term{local-duality}{局部对偶}定理\footnote{参见 SGA 5 I。}。在\Term{algebra}{代数}几何中，写下这些文字时，只有最后一个尚未证明（不过在\Term{characteristic}{特征} $0$ 时，使用 Hironaka 的\Term{resolution-of-singularities}{奇点消解}已可证明）。此外，在超越框架中，现在已经有 Verdier 最近证明的整体与\Term{local-duality}{局部对偶}性。\footnote{编者注（20）：见 Verdier J.-L., “Dualité dans la cohomologie des espaces localement compacts”, 载 Séminaire Bourbaki, 第 9 卷，Société mathématique de France, Paris, 1995，第 300 讲，337--349 页。}这里只指出，在问题 4.2 与 4.3 的\Term{homology}{同调}版本（现在已值得称为猜想）的陈述中，“无穷远处”的条件 a) 与 c) 肯定是多余的，重要的只是
\SourcePage{162}
$X-Y$ 的\Term{local-cohomology}{局部上同调}结构，例如可假设它\Term{local-scheme}{局部}为\Term{dimension}{维数} $\geq n$ 的\Term{complete-intersection}{完全交}。另外，比如在 4.3 中，$Y$ 是\Term{hyperplane-section}{超平面截面}这一事实应当不起\Term{action}{作用}，而应当能替换为仅要求 $X$ 紧且 $X-Y$ 为 Stein（即在\Term{algebra}{代数}几何情形，$X$ 在 $k$ 上\Term{proper}{固有}且 $X-Y$ \Term{affine}{仿射}；如前所说，这个猜想的\Term{homology}{同调}版本已经对复数域 $\mathbf C$ 上的\Term{algebraic-space}{代数空间}证明）。\footnote{参见第 XIV 讲。}

在 $\pi_i^x(X)$ 的定义（原书第 146 页）中，应当假设 $i\geq2$。对于 $i=0,1$，没有合理的 $\pi_i^x(X)$ 的定义，应当把它们换成 $H_0^x(X)$ 和 $H_1^x(X)$，分别定义为自然\Term{homomorphism}{同态}
\[
\varprojlim H_0(U-(x),\mathbf Z)\longrightarrow\varprojlim H_0(U,\mathbf Z)
\]
的\Term{cokernel}{余核}和核。

严格说来，为了表述方便，可以在 $i\leq1$ 时令 $\pi_i^x(X)=H_i^x(X)$；否则，就必须用对应的关于 $H_0^x$、$H_1^x$ 的命题，补足后面有关 $\pi_i^x$ 的命题。如果 $x$ 是 $X$ 的\Term{isolated-point}{孤立点}，宜令 $i\ne0$ 时 $\pi_i^x(X)=0$，而 $\pi_0^x(X)=H_0^x(X)=\mathbf Z$。

关于各个 $\pi_i^x(U,f)$ 彼此\Term{isomorphism}{同构}的断言，只有当 $x$ 在其邻域内不使 $X$ 不\Term{connected}{连通}，即 $i=0,1$ 时 $\pi_i^x(X)=0$ 才正确。在一般情形，$\pi_i^x(X)$ 只能表示一族未必彼此\Term{isomorphism}{同构}的群；不过，写成 $\pi_i^x(X)=0$ 仍有显然的意义。

在原书第 146 页，我预期对 $x\in Y$、$i\leq n$，\Term{local-scheme}{局部}\Term{homotopy}{同伦}\Term{invariants}{不变量} $\pi_i^x(X)$ 的消失，应当蕴含 $\pi_i(X-Y)\to\pi_i(X)$ 在 $i<n-d$ 时\Term{bijection}{双射}、在 $i=n-d$ 时\Term{surjectivity}{满射}；对此应当谨慎，因为在这里的背景中，我们不像在\Term{algebra}{代数}几何中那样，拥有也必须满足这些\Term{local-scheme}{局部}条件的“一般”点。为此，无疑需要借助相对\Term{local-scheme}{局部}\Term{homotopy}{同伦}\Term{invariants}{不变量}
\[
\pi_i^Y(X,f)=\pi_i^Y(X,x)
 =\varprojlim_U\pi_{i-1}(U-U\cap Y,f(t))\qquad\text{当 }i\geq2
\]
（而在 $i=0,1$ 时如上作专门定义），其中 $Y$ 是 $X$ 的\Term{closed-subset}{闭子集}；或者，把关于 $X$ 的假设表述为 $X$ 的\Term{local-ring}{局部环}的谱的\Term{topology}{拓扑}性质（对于\Term{etale-topology}{平展拓扑}），以补偿\Term{generic-point}{一般点}的缺失，这样就能重新获得\Term{generic-point}{一般点}。相同的保留意见也适用于将猜想 4.2 和 4.3 推广到不假设 $X-Y$ \Term{local-scheme}{局部}为\Term{complete-intersection}{完全交}的情形；这种推广已由这些猜想中条件 b) 的陈述所提示。

为了表述上述结果所提示的、经过订正的猜想 4.2 与 4.3，应当给出如下定义。

\SourcePage{163}
\begin{SGAresult}{定义}{1}
设 $X$ 是\Term{topological-space}{拓扑空间}，$Y$ 是 $X$ 的\Term{locally-closed}{局部闭}\Term{subset}{子集}，$n$ 为整数。若对每个 $x\in Y$，在 $i<n$ 时都有 $\pi_i^Y(X,x)=0$，则称 $X$ 沿 $Y$ 的\Term{homotopical-depth}{同伦深度} $\geq n$，记作 $\operatorname{prof}\operatorname{hpt}_Y(X)\geq n$。
\end{SGAresult}
这应当\Term{equivalence}{等价}于说：对于 $X$ 的每个\Term{open-set}{开集} $X'$，以及每个 $x\in X'\cap U=U'$（其中 $U=X-Y$），\Term{canonical}{典范}\Term{homomorphism}{同态}
\[
\pi_i(U',x)\longrightarrow\pi_i(X',x)
\]
在 $i<n-1$ 时为\Term{isomorphism}{同构}，在 $i=n-1$ 时为\Term{monomorphism}{单态射}。\footnote{编者注（21）：若偶对 $(X,Y)$ 还是多面体偶对，这一\Term{equivalence}{等价}成立；参见 Eyral C., “Profondeur homotopique et conjecture de Grothendieck”, Ann. Sci. Éc. Norm. Sup. (4) 33 (2000), 第 6 期，823--836 页。}

\begin{SGAresult}{定义}{2}
设 $X$ 为\Term{complex-analytic-space}{复解析空间}，$n$ 为整数；若对 $X$ 的每个\Term{locally-closed}{局部闭}\Term{analytic-subset}{解析子集} $Y$，都有
\[
\operatorname{prof}\operatorname{htp}_Y(X)\geq n-\dim Y,\tag{$x$}
\]
则称 $X$ 的\Term{rectified-homotopical-depth}{校正同伦深度}为 $n$。\footnote{这些笔记的第一版中，我们使用了“真正\Term{homotopical-depth}{同伦深度}”一词。在这一版中，我们遵循 EGA IV 10.8.1。}
（其中，当然，$\dim Y$ 指 $Y$ 的复\Term{dimension}{维数}。）
\end{SGAresult}
这应当\Term{equivalence}{等价}于说：对于 $X$ 的每个\Term{locally-closed}{局部闭}\Term{irreducible}{不可约}\Term{analytic-subset}{解析子集} $Y$，存在 $Y$ 的一个\Term{dimension}{维数} $<\dim Y$ 的\Term{closed-analytic-subset}{闭解析子集} $Z$，使得把 $Y$ 换成 $Y-Z$ 时，关系 ($x$) 成立。例如，这样一来，定义 2 中就能只考虑 $Y$ \Term{nonsingular}{非奇异}的情形。\footnote{编者注（22）：后面的所有猜想，经过适当校正——如果允许我这样说的话——都因 Hamm 和 Lê Dũng Tráng 的工作而成为定理，见 Hamm H.A. \& Lê Dũng Tráng, “Rectified homotopical depth and Grothendieck conjectures”, 载 The Grothendieck Festschrift, Vol. II, Progr. Math., 第 87 卷，Birkhäuser, Boston, 1990，311--351 页，以下引用为 [HL]。至于校正\Term{depth}{深度}的两个被猜测为\Term{equivalence}{等价}的定义，它们甚至还\Term{equivalence}{等价}于第三个用 Withney 分\Term{sheaf}{层}表述的定义（参见所引文献，定理 1.4）。}
下面的猜想是纯\Term{topology}{拓扑}性质的，具有“\Term{local-scheme}{局部} Hurewicz 定理”的性质。

\begin{SGAresult}{猜想}{A（“\Term{local-scheme}{局部} Hurewicz 定理”）}
\footnote{编者注（23）：如 [HL] 例 3.1.3 所指出的，这个猜想对 $X=\{z\in\mathbf C^n\mid z_1^2+z_2^3+\cdots+z_n^3=0\}$、$n\geq4$，且 $Y$ 只含原点时就已经是错误的。但经适当修改后，它成立（所引文献定理 3.1.4）。}
设 $X$ 为\Term{topological-space}{拓扑空间}，$Y$ 为\Term{locally-closed}{局部闭}\Term{subset}{子集}，必要时施加某些“\Term{smooth}{光滑}性”条件，例如偶对 $(X,Y)$ \Term{local-scheme}{局部}可三角剖分；$n$ 为整数，$n\geq3$。要有 $\operatorname{prof}\operatorname{htp}_Y(X)\geq n$，充要条件是
\[
\underline H^i_Y(\mathbf Z_X)=0\qquad\text{当 }i<n
\]
\SourcePage{164}
（此时称 $X$ 沿 $Y$ 的\Term{cohomological-depth}{上同调深度} $\geq n$），并且\Term{local-scheme}{局部}\Term{fundamental-group}{基本群}
\[
\pi_2^Y(X,x)=\varprojlim_{U\ni x}\pi_1(U-U\cap Y)
\]
都为零（此时称 $X$ 沿 $Y$ “纯”）。
\end{SGAresult}

注意，若 $X$ 是\Term{analytic-space}{解析空间}，$Y$ 是\Term{analytic-subspace}{解析子空间}，而且 $X$ 沿 $Y$ 纯，则对每个 $x\in Y$，\Term{local-ring}{局部环} $\mathcal O_{X,x}$ 以及它对那些包含定义 $x$ 处 $Y$ 的芽的\Term{ideal}{理想}的\Term{prime-ideal}{素理想}的\Term{localization}{局部化}（即在 $\operatorname{Spec}(\mathcal O_{X,x})=X_x\to X$ 下 $Y$ 的\Term{inverse-image}{逆像} $Y_x$ 内），都在第 X 讲的意义下为纯；其逆命题似乎也很可能正确。对\Term{cohomological-depth}{上同调深度}也有类似注记，不过要明确，我们在 $\operatorname{Spec}(\mathcal O_{X,x})$ 上使用\Term{etale-topology}{平展拓扑}。

于是猜想 4.1 推广为如下猜想。
\begin{SGAresult}{猜想}{B（“\Term{purity}{纯性}”）}
\footnote{编者注（24）：[HL] 已证明这个猜想，即使 $E$ \Term{singular}{奇异}也成立：这就是定理 3.2.1。}
设 $E$ 为\Term{analytic-space}{解析空间}，$X$ 为 $E$ 的\Term{analytic-subset}{解析子集}。假设 $E$ 在 $x\in X$ 处\Term{nonsingular}{非奇异}且\Term{dimension}{维数}为 $N$，而 $X$ 在每个点附近可以用 $p$ 个解析方程描述。则 $X$ 的\Term{rectified-homotopical-depth}{校正同伦深度} $\geq N-p$。
\end{SGAresult}
特别地，在每个点处\Term{dimension}{维数}为 $n$ 的\Term{local-complete-intersection}{局部完全交}，其\Term{rectified-homotopical-depth}{校正同伦深度}应当 $\geq n$；这无非就是猜想 4.1。

猜想 4.2 和 4.3 分别推广为：
\begin{SGAresult}{猜想}{C（“\Term{local-scheme}{局部} Lefschetz”）}
\footnote{编者注（25）：[HL] 已证明这个猜想，甚至包括后面注记中的强形式，参见所引文献定理 3.3.1。}
设 $X$ 为\Term{analytic-space}{解析空间}，$Y$ 为\Term{closed-analytic-subset}{闭解析子集}，$x$ 为 $Y$ 的一点。假设 $X-Y$ 在 $x$ 的邻近处为 Stein（例如 $Y$ 在 $x$ 处由一个方程定义），并且 $X-Y$ 在 $x$ 的邻近处\Term{rectified-homotopical-depth}{校正同伦深度} $\geq n$（例如在 $X-Y$ 中邻近 $x$ 的每一点处，都是\Term{dimension}{维数} $\geq n$ 的\Term{complete-intersection}{完全交}，参见猜想 B）。则\Term{canonical}{典范}\Term{homomorphism}{同态}
\[
\pi_i^x(Y)\longrightarrow\pi_i^x(X)
\]
在 $i<n-1$ 时为\Term{isomorphism}{同构}，在 $i=n-1$ 时为\Term{epimorphism}{满态射}。
\end{SGAresult}

\begin{SGAresult}{猜想}{D（“整体 Lefschetz”）}
\footnote{编者注（26）：[HL] 同样已证明这个猜想，甚至包括后面注记中的强形式，参见所引文献定理 3.4.1。}
设 $X$ 为紧\Term{analytic-space}{解析空间}，$Y$ 为 $X$ 的\Term{analytic-subspace}{解析子空间}，使得 $U=X-Y$ 为 Stein，并且\Term{rectified-homotopical-depth}{校正同伦深度}
\SourcePage{165}
$\geq n$（例如在每个点处都是\Term{dimension}{维数} $\geq n$ 的\Term{complete-intersection}{完全交}）。则\Term{canonical}{典范}\Term{homomorphism}{同态}
\[
\pi_i(Y)\longrightarrow\pi_i(X)
\]
在 $i<n-1$ 时为\Term{isomorphism}{同构}，在 $i=n-1$ 时为\Term{epimorphism}{满态射}。
\end{SGAresult}

\textbf{注记。}在命题 C 和 D 中，若将 $X-Y$ 为 Stein 的假设换成 $X-Y$ 是 $c+1$ 个 Stein \Term{open-set}{开集}的并（它将起到\Term{topology}{拓扑}“凹性”假设的\Term{action}{作用}），则只须在结论中把 $n$ 换成 $n-c$。\footnote{编者注（27）：最后指出 Fulton 的下述结果，可与原书第 127 页编者注（4）中引用的 Fulton--Hansen 结果比较：设 $X$、$H$ 为 $\mathbf P^m_{\mathbf C}$ 的\Term{closed-subscheme}{闭子概形}，$n$ 为 $X$ 的\Term{dimension}{维数}，$d$ 为 $H$ 的\Term{codimension}{余维数}。则映射
\[
\pi_i(X,X\cap H)\longrightarrow\pi_i(\mathbf P^n_{\mathbf C},H)
\]
在 $i\leq n-d$ 时为\Term{isomorphism}{同构}，在 $i=n-d-1$ 时\Term{surjectivity}{满射}；见 Fulton W., “Connectivity and its applications in algebraic geometry”, 载 Algebraic geometry (Chicago, Ill., 1980), Lect. Notes in Math., 第 862 卷，Springer, Berlin--New York, 1981，26--92 页。}

最后，在“整体情形”D 中，明确写出有关\Term{fundamental-group}{基本群}的猜想（令 $n=3$ 所得）：

\begin{SGAresult}{猜想}{D'（\Term{fundamental-group}{基本群}的整体 Lefschetz 定理）}
\footnote{编者注（29）：[HL] 已证明这个猜想，参见所引文献定理 3.5.1。}
设 $X$ 为复数域上的紧\Term{analytic-space}{解析空间}，$Y$ 为\Term{closed-analytic-subset}{闭解析子集}，使得 $U=X-Y$ 为 Stein。另外假设满足下列条件：

(i) 对每个 $x\in U$，\Term{local-scheme}{局部}\Term{fundamental-group}{基本群} $\pi_2^x(X,x)$ 为零（即 $X$ “在 $x$ 处纯”），或者只要求\Term{local-ring}{局部环} $\mathcal O_{X,x}$ 纯。

(ii) $U$ 的各点的\Term{local-ring}{局部环}“在\Term{dimension}{维数} $\geq2$ 上\Term{connected}{连通}”。

(iii) $U$ 的各点的\Term{local-ring}{局部环}的\Term{dimension}{维数}均 $\geq3$。

在这些条件下，对每个 $x\in Y$，\Term{homomorphism}{同态}
\[
\pi_1(Y,x)\longrightarrow\pi_1(X,x)
\]
为\Term{isomorphism}{同构}（而 $\pi_2(Y,x)\to\pi_2(X,x)$ 为\Term{epimorphism}{满态射}）。
\end{SGAresult}

注意，若 $U$ \Term{local-scheme}{局部}为\Term{dimension}{维数} $\geq3$ 的\Term{complete-intersection}{完全交}，则 $U$ 上的\Term{local-scheme}{局部}条件 (i)、(ii)、(iii) 都满足。从\Term{algebra}{代数}几何观点看（当 $U$ 来自一个仍记为 $U$ 的\Term{scheme}{概形}时），条件 (i) 至 (iii) 对应于关于\Term{local-scheme}{局部}\Term{invariants}{不变量} $\pi_i^x(U)$ 的假设，即分别对满足 $\operatorname{trdeg}k(x)/k=0,1,2$ 的点 $x$，当 $i<3-\operatorname{trdeg}k(x)/k$ 时，$\pi_i^x(U)=0$。若 $X$ \Term{projective}{射影}而 $Y$ 为\Term{hyperplane-section}{超平面截面}，则 $U$ 上的整体条件（$U$ 为 Stein）得到满足。

\SourcePage{166}
\SGAsection{}{参考文献}
\begin{enumerate}
\item[{[1]}] S. Abhyankar，“Local uniformisation on algebraic surfaces over ground fields of characteristics $p\ne0$”（\Term{characteristic}{特征} $p\ne0$ \Term{base-field}{基域}上\Term{algebra}{代数}曲面的\Term{local-scheme}{局部}一致化），Annals of Math. 63 (1956)，491--526 页。
\item[{[2]}] W. L. Chow，“On the theorem of Bertini for local domains”（论\Term{local-scheme}{局部}\Term{integral-domain}{整环}的 Bertini 定理），Proc. Nat. Acad. Sci. U.S.A. 44 (1958)，580--584 页。
\item[{[3]}] M. Greenberg，“Schemata over local rings”（\Term{local-ring}{局部环}上的\Term{scheme}{概形}），Annals of Math. 73 (1961)，624--648 页。
\item[{[4]}] M. Kneser，“Über die Darstellung algebraischer Raumkurven als Durchschnitte von Flächen”（论以曲面之交表示\Term{algebraic-space}{代数空间}曲线），Archiv der Math. XI (1960)，157--158 页。
\item[{[5]}] D. Mumford，“The topology of normal singularities of an algebraic surface, and a criterion for simplicity”（\Term{algebra}{代数}曲面\Term{normal}{正规}奇点的\Term{topology}{拓扑}及单性判据），Publ. Math. Inst. Hautes Études Sci. 9 (1961)，5--22 页。
\item[{[6]}] A. Néron，“Modèles minimaux des variétés abéliennes sur les corps locaux et globaux”（\Term{local-scheme}{局部}域和整体域上\Term{abelian-variety}{阿贝尔簇}的\Term{minimal-model}{极小模型}），Publ. Math. Inst. Hautes Études Sci. 21 (1964)，5--128 页。
\item[{[7]}] A. H. Wallace，Homology Theory of algebraic Varieties（\Term{algebraic-variety}{代数簇}的\Term{homology}{同调}理论），Pergamon Press，1958。
\item[{[8]}] S. Abhyankar，“Resolution of singularities of arithmetical surfaces”（算术曲面的\Term{resolution-of-singularities}{奇点消解}），载 Arithmetical Algebraic Geometry，Harper and Row, New-York，1965，111--152 页。
\item[{[HL]}] H. A. Hamm \& Lê Dũng Tráng，“Rectified homotopical depth and Grothendieck conjectures”（\Term{rectified-homotopical-depth}{校正同伦深度}与 Grothendieck 猜想），载 The Grothendieck Festschrift, Vol. II, Progr. Math.，第 87 卷，Birkhaüser, Boston，1990，311--351 页。
\end{enumerate}
